% !TeX encoding = utf8
% !TeX spellcheck = pl-PL

\documentclass[thesis=mgr,faculty=ee]{EE-dyplom} 



\instytut{Instytut Sterowania i Elektroniki Przemysłowej}
\kierunek{Informatyka Stosowana}
\specjalnosc{Informatyka w Biznesie}
\title{System śledzenia i reidentyfikacji obiektów
w nagraniach z wielu kamer monitoringu z
wykorzystaniem mapy terenu}
\engtitle{Object Tracking and Re-Identification System
for Multi-Camera Surveillance Footage Using
a Map of the Monitored Area}
\album{340803}
\author{Natalia Turska}
\promotor{dr inż. Witold Czajewski}
\date{\the\year{}}
\miejscowosc{Warszawa}

\streszczeniepracy{
Streszczenie pracy powinno w zwięzły sposób opisywać to czego dotyczy praca, co jest jej celem, jakie przyjęto założenia, co w ramach pracy zrobiono, co przebadano, jakie rozwiązania zaproponowano i jakie wyniki osiągnięto. Jeśli coś sprawiło jakiś problem, można to także ująć w streszczeniu i skomentować czy udało się ów problem rozwiązać i jak, czy też nie (nic w tym złego, jeśli czegoś się nie udało do końca zrobić). 
\\
Streszczenie nie powinno być przeładowane, powinno mieć około 200 słów i zawierać najważniejsze informacje na temat pracy i najważniejsze osiągnięcia, nie należy więc podawać szczegółowych wyników czy też opisywać struktury pracy - na to miejsce znajduje się w samej pracy.
\\
Streszczenie zwykle pisze się na samym końcu pracy - z oczywistych względów.

\noindent\textit{W niniejszej pracy magisterskiej zaprezentowane zostały próby rozwiązania problemu klasyfikacji poziomu zabrudzenia chodników z wykorzystaniem głębokich sieci neuronowych. Wykorzystując dostępny zbiór zdjęć sklasyfikowanych do sześciu klas poziomu zanieczyszczeń, wygenerowano zrównoważone zbiory: sześcioklasowy oraz trzyklasowy. Następnie wybrano trzy sieci różnego rodzaju, aby sprawdzić ich efektywność w rozróżnianiu poszczególnych, często bardzo zbliżonych do siebie zdjęć. Sieć TResNet została wybrana spośród wiodących sieci w rankingach klasyfikacyjnych. Sieć ShuffleNet V2 została wybrana spośród dostępnych implementacji sieci w ramach popularnej biblioteki OpenMMLab. Sieć PMG została wybrana spośród sieci przygotowanych do rozwiązywania specyficznych problemów klasyfikacyjnych dla obrazów o niewielkich różnicach. Wymienione sieci wytrenowano na wygenerowanych zbiorach i osiągnięto dokładność rzędu 80\% i 90\%, odpowiednio dla problemu sześcio-~i trzyklasowego. Najlepszą siecią, w każdym przypadku, okazała się sieć PMG, a wyniki pozostałych sieci były zbliżone.}
} % koniec streszczenia
%%%%%%%%%%%%%%%%%%%%%%%%%%%%%%%%%%%%%%%%%%%%%%%%%%%
\slowakluczowe{istotne, słowa, kluczowe} 
%%%%%%%%%%%%%%%%%%%%%%%%%%%%%%%%%%%%%%%%%%%%%%%%%%%
\thesisabstract{
% Należy zadbać oto, by tłumaczenie było poprawne i zgodne z oryginałem odnośnie meritum. Koniecznie trzeba sprawdzić tłumaczenie zrobione przez SI, która ostatnio bardzo dobrze tłumaczy, ale czasami może czegoś nie zrozumieć.
\textit{This engineering thesis presents an attempt to solve the problem of sidewalk dirt level classification using deep neural networks. Using an available set of images classified into six classes of dirt levels, balanced sets of six classes and three classes were generated. Three networks of different types were then selected to test their effectiveness in discriminating between individual, often very similar, images. TResNet was selected among the leading networks in the classification rankings. The ShuffleNet V2 network was selected from available network implementations within the popular OpenMMLab library. The PMG network was selected from networks prepared to solve specific classification problems for images with small differences. The aforementioned networks were trained on the generated sets and an accuracy of 80\% and 90\% was achieved for the six-class and three-class problem, respectively. The PMG network proved to be the best network, in each case, and the results of the other networks were similar.}
} 
\thesiskeywords{keywords, that, are, indicative}  

\begin{document}
\frontpages

\chapter{Wstęp}\label{chap:wprowadzenie}

\chapter{Stan wiedzy}\label{chap:stan_wiedzy}
Najwcześniejsze systemy telewizji przemysłowej \textit{(ang. Closed Circuit Television -- CCTV)}
powstały w pierwszej połowie XX wieku i umożliwiały jedynie bieżącą obserwację zdarzeń, ze względu
na brack technologii pozwalających na rejestrowanie i przechowywanie sygnału wideo. Rozwój nośników
szpulowych umożliwił rejestrowanie materiału z monitoringu. Systemy te wymagały były jednak
uciążliwe w obsłudze - wymagały ręcznej wymiany taśm magnetycznych oraz ręcznego nawijania taśmy na
szpulę odbiorczą. Wraz z postępem technologicznym metody zapisu ulegały systematycznej poprawie, co
przyczyniło się do popularyzacji systemów bezpieczeństwa. Już w latach 90. XX opracowano
technologię multipleksowania cyfrowego, umożliwiającą jednoczesne nagrywanie z kilku kamer, a także
nagrywanie w trybie poklatkowym i zapis wyłącznie przy wykryciu ruchu~\cite{kreugle}.

Istotny przełom w rozwoju systemów wideonadzoru przyniosła wizja komputerowa, rozwijana od lat 70.
XX wieku, której nadrzędnym celem stało się pełne zautomatyzowane zrozumienie analizowanej sceny.
Prawdziwą rewolucję w tej dziedzinie umożliwiło jednak dopiero znaczne zwiększenie mocy
obliczeniowej, rozwój algorytmów uczenia maszynowego oraz powstanie wielkich, oznaczonych zbiorów
danych (m.in. ImageNet, Microsoft COCO czy LVIS). Bazy te dostarczyły niezbędnych informacji do
trenowania zaawansowanych modeli oraz pozwoliły na obiektywne mierzenie postępów w dziedzinie
rozpoznawania i segmentacji obrazu~\cite{szeliski2022}.

%%%%%%%%%%%%%%%%%%%%%%%%%%%%%%%%%%%%%%%%%%%%%%%%%%%%%%%%%%%
%%%%%%%%%%%%%%%%%%%%%%%%%%%%%%%%%%%%%%%%%%%%%%%%%%%%%%%%%%%

\section{Analiza obrazu i podstawowe definicje}

\subsection{Komputerowa reprezentacja obrazu}
Komputer, analizując obraz, sprowadza go do postaci macierzy pikseli, którym przypisane są wartości
liczbowe zależne od koloru. Kształt tej macierzy zależy od przestrzeni barw, w jakiej
reprezentowany jest dany obraz. W przypadku skali szarości obraz jest przedstawiany w postaci
macierzy dwuwymiarowej, której elementy przyjmują wartości od 0 do 255, gdzie 0 oznacza czerń, a
255 biel~\cite{gonzalezwoods}.

Obrazy przedstawione w pełnym spektrum kolorów zapisywane są najczęściej w przestrzeni RGB
\textit{(ang. Red, Green, Blue)} lub HSV \textit{(ang. Hue, Saturation, Value)}. W takim przypadku
obraz opisywany jest za pomocą trzech macierzy, a każdemu pikselowi przypisywane są trzy wartości z
zakresu od 0 do 255. W przestrzeni RGB, która jest wykorzystywana częściej, wartości te odpowiadają
odpowiednio natężeniu koloru czerwonego, zielonego i niebieskiego~\cite{gonzalezwoods}. Takie
przedstawienie pozwala na odwzorowanie pełnego spektrum barw w formie dogodnej do przetwarzania
komputerowego.

Liczba pikseli tworzących obraz jest równoznaczna z jego rozdzielczością~\cite{gonzalezwoods}, a
więc bezpośrednio wpływa na poziom szczegółowości widocznych obiektów i możliwości ich poprawnej
detekcji. Nagrania wideo składają się z pojedynczych obrazów nazywanych w tym kontekście klatką lub
ramką \textit{(ang. frame)}. Ciąg kolejnych klatek wyświetlanych w krótkich odstępach czasu tworzy
wrażenie ruchu i stanowi strumień wideo. Liczba klatek na sekundę, wpływa na płynność nagrania oraz
możliwość dokładnej analizy ruchu obiektów w czasie.

\subsection{Model kamery i rzutowanie na plan terenu}
Dla poprawnej analizy nagrań wideo należy wziąć pod uwagę sposób ich powstawania w kamerze. Aby
możliwe było poprawne wyznaczanie położenia obiektów w scenie, stosuje się proces kalibracji
kamery, który pozwala określić zależność między współrzędnymi punktów w przestrzeni rzeczywistej a
ich położeniem w obrazie. Podstawą takiego opisu jest model otworkowy kamery \textit{(ang. pinhole
  camera model)}, w którym obraz powstaje jako rzut perspektywiczny punktów trójwymiarowej sceny na
dwuwymiarową płaszczyznę obrazu~\cite{hartleyzisserman}.

Pełny model kamery można zapisać w postaci iloczynu macierzy wewnętrznej oraz macierzy zewnętrznej.
Macierz wewnętrzna opisuje właściwości samej kamery i sposób odwzorowania punktów w układzie kamery
na współrzędne obrazu. Zawiera ona między innymi ogniskowe w osiach poziomej i pionowej, położenie
punktu głównego oraz ewentualny parametr skośności osi. Parametry te określają więc geometryczne
cechy obrazowania i wpływają na skalę oraz położenie obrazu w układzie pikseli~\cite{szeliski2022}.

Grupa parametrów zewnętrznych, opisuje położenie i orientację kamery względem układu współrzędnych
sceny. Obejmują one macierz obrotu oraz wektor translacji. Parametry te opisują transformację
punktów ze współrzędnych świata do układu współrzędnych kamery. Dzięki temu możliwe jest
określenie, z którego miejsca kamera obserwuje scenę oraz w jakim kierunku jest
skierowana~\cite{szeliski2022}.

Złożenie macierzy wewnętrznej i zewnętrznej daje pełną macierz projekcji kamery \(\mathbf{P} =
\mathbf{K}[\mathbf{R} \mid \mathbf{t}]\), która opisuje kompletną transformację punktu z
trójwymiarowej przestrzeni sceny do dwuwymiarowego układu pikselowego
obrazu~\cite{hartleyzisserman}.

W praktyce rzeczywisty obraz może odbiegać od idealnego modelu z powodu dystorsji soczewki, która
powoduje zniekształcenia, są one szczególnie widoczne na obrzeżach obrazu i wynikają z właściwości
układu optycznego. Z tego względu w procesie kalibracji często wyznacza się również współczynniki
pozwalające na korekcję tych zniekształceń przed dalszą analizą obrazu~\cite{szeliski2022}.

Do analizy wielu obrazów tej samej powierzchni lub terenu służy analiza homografii. W dziedzinie
widzenia komputerowego dowolne dwa obrazy tej samej płaskiej powierzchni w przestrzeni są powiązane
przekształceniem rzutowym, zwanym właśnie homografią. Opisuje się ją za pomocą macierzy
\(\mathbf{H}\) o wymiarach \(3 \times 3\), która odwzorowuje punkty z jednej płaszczyzny na drugą z
uwzględnieniem zniekształceń perspektywy. Macierz ta działa na punkty zapisane we współrzędnych
jednorodnych i przyjmuje postać:
\[
  \lambda
  \begin{bmatrix}
    x' \\
    y' \\
    1
  \end{bmatrix}
  =
  \mathbf{H}
  \begin{bmatrix}
    x \\
    y \\
    1
  \end{bmatrix},
\]
gdzie \((x, y)\) oznacza współrzędne punktu w pierwszym układzie, \((x', y')\) współrzędne punktu
po przekształceniu, natomiast \(\lambda\) jest współczynnikiem skali wynikającym z użycia
współrzędnych jednorodnych. Taki zapis pozwala ująć w jednej macierzy przesunięcie, obrót, zmianę
skali oraz zniekształcenie perspektywiczne~\cite{zhang2000, szeliski2022}.

W kontekście systemów wielokamerowych rzutowanie na plan terenu umożliwia przedstawienie
trajektorii obiektów w jednej, wspólnej przestrzeni geometrycznej, niezależnie od położenia i
orientacji poszczególnych kamer.

Należy jednak zaznaczyć, że poprawność takiego rzutowania zależy od jakości kalibracji kamery oraz
od tego, w jakim stopniu obserwowana scena spełnia założenie o płaskości terenu. W przypadku
nierównego podłoża lub obiektów o znacznej wysokości konieczne jest zastosowanie bardziej
rozbudowanych modeli geometrycznych.

\subsection{Detekcja ruchu i odejmowanie tła}

Technika odejmowania tła \textit{(ang. background subtraction)} umożliwia rozpoznanie ruchomych
obiektów na nagraniach wideo dla statycznych scen poprzez oddzielenie pierwszego planu od tła.
Przyjmując założenie, że tło pozostaje wzglęgnie niezmienne (mogą pojawiać się zmiany oświetlenia,
zabrudzenia itp.), możliwe jest opisanie tła za pomocą modelu statystyczngo. Stworzony model tła
odejmowany jest od obrazu, pozostawiając jedynie anomalie, czyli obiekty znajdujące się w ruchu,
pierwszy plan zwracany jest zazwyczaj w postaci maski.~\cite{zivkovic2004}

\begin{figure}[H]
  \centering
  \includegraphics[width=.8\textwidth]{background-substraction.jpg}
  %\vspace*{-10mm}
  \caption{Wizualizacja procesu odejmowania tła. Źródło:~\cite{background_subtraction}}
  \label{fig:mog2}
\end{figure}

W kontekście projektowania zaawansowanych systemów śledzenia wielu obiektów i reidentyfikacji,
algorytmy odejmowania tła pozwalają na ograniczenie obszaru poszukiwań, system może przetwarzać
wyłącznie wykadrowane obszary ruchu (ang. Regions of Interest – ROI).

\subsubsection{MOG2 (Mixture of Gaussians 2)}
Algorytm MOG2 stanowi bezpośrednie rozwinięcie metody mieszanin rozkładów Gaussa zaproponowanej
przez Stauffera i Grimsona. Model GMM \textit{(ang.~Gaussian Mixture Model)} zakłada stałą liczbę
komponentów gaussowskich dla każdego piksela w scenie, co może prowadzić do nadmiernej złożoności
obliczeniowej dla jednorodnych i statycznych obszarów obrazu. MOG2, opracowany przez Zivkovica,
wprowadza adaptacyjny dobór liczby rozkładów dla każdego piksela z osobna~\cite{zivkovic2004}.

W~algorytmie MOG2 prawdopodobieństwo zaobserwowania określonej wartości piksela $\vec{x}$ w~chwili
$t$ w~wielowymiarowej przestrzeni barw (np. RGB) jest opisywane przez ważoną sumę $M$ rozkładów
normalnych~\cite{stauffer1999, zivkovic2004}:

\begin{equation}
  \hat{p}(\vec{x} \mid X_T, \text{BG}+\text{FG}) = \sum_{m=1}^{M} \hat{\pi}_m \mathcal{N}(\vec{x}; \hat{\vec{\mu}}_m, \hat{\sigma}_m^2 \mathbf{I})
  \label{eq:mog2_prob}
\end{equation}

gdzie $\hat{\pi}_m$ oznacza oszacowaną wagę $m$-tego komponentu, $\hat{\vec{\mu}}_m$ jest jego
wektorem średnich, natomiast $\hat{\sigma}_m^2 \mathbf{I}$ reprezentuje macierz
kowariancji~\cite{stauffer1999, zivkovic2004}. Dla uproszczenia obliczeń zakłada się, że
poszczególne kanały kolorów są od siebie niezależne i~posiadają tę samą wariancję, co sprowadza
macierz kowariancji do iloczynu wariancji $\hat{\sigma}_m^2$ oraz macierzy jednostkowej
$\mathbf{I}$~\cite{stauffer1999, zivkovic2004}. Wagi mieszaniny są nieujemne i~podlegają
normalizacji, spełniając warunek $\sum_{m=1}^{M} \hat{\pi}_m = 1$~\cite{stauffer1999,
  zivkovic2004}.

Dla każdej nowej klatki obrazu i~każdej nowo zaobserwowanej wartości piksela $\vec{x}(t)$
realizowana jest rekurencyjna aktualizacja parametrów modelu~\cite{stauffer1999, zivkovic2004}.
Najpierw następuje porządkowanie komponentów według stosunku wagi do odchylenia standardowego
($\hat{\pi}_m / \hat{\sigma}_m$), co faworyzuje rozkłady o~wysokiej wadze i~niskiej wariancji jako
najbardziej prawdopodobne definicje tła~\cite{stauffer1999}. Dopasowanie próbki do istniejących
komponentów określa się na podstawie odległości Mahalanobisa~\cite{stauffer1999, zivkovic2004}.
Jeśli kwadrat tej odległości jest mniejszy od przyjętego progu (odpowiadającego zazwyczaj trzem
odchyleniom standardowym), komponent uznaje się za dopasowany~\cite{stauffer1999}. Zmienna
przynależności $o_m^{(t)}$ przyjmuje wówczas wartość 1 dla najlepiej dopasowanego rozkładu i~0 dla
pozostałych~\cite{stauffer1999}. Aktualizacja parametrów zachodzi według następujących równań
rekurencyjnych~\cite{stauffer1999, zivkovic2004}:

\begin{equation}
  \hat{\pi}_m \leftarrow \hat{\pi}_m + \alpha (o_m^{(t)} - \hat{\pi}_m)
  \label{eq:mog2_weight_update}
\end{equation}
\begin{equation}
  \hat{\vec{\mu}}_m \leftarrow \hat{\vec{\mu}}_m + o_m^{(t)} \left( \frac{\alpha}{\hat{\pi}_m} \right) \vec{\delta}_m
  \label{eq:mog2_mean_update}
\end{equation}

\begin{equation}
  \hat{\sigma}^2_m \leftarrow \hat{\sigma}^2_m + o_m^{(t)} \left( \frac{\alpha}{\hat{\pi}_m} \right) \left( \vec{\delta}_m^T \vec{\delta}_m - \hat{\sigma}^2_m \right)
  \label{eq:mog2_var_update}
\end{equation}

gdzie $\vec{\delta}_m = \vec{x}(t) - \hat{\vec{\mu}}_m$, natomiast $\alpha$ określa współczynnik
decydujący o tempie zmian parametrów rozkładu, $\alpha = 1/T$, gdzie $T$ jest rozważanym interwałem
czasowym (liczbą klatek)~\cite{stauffer1999, zivkovic2006}. Jeśli dla danego piksela nie zostanie
znaleziony żaden dopasowany komponent, tworzony jest nowy rozkład o~parametrach: $\hat{\pi}_{M+1} =
  \alpha$, $\hat{\vec{\mu}}_{M+1} = \vec{x}(t)$ oraz wariancji początkowej $\hat{\sigma}_{M+1}^2 =
  \sigma_0^2$~\cite{stauffer1999}. Jeśli liczba komponentów przekroczy dopuszczalne maksimum, rozkład
o~najmniejszej wadze jest usuwany~\cite{stauffer1999}.

\subsubsection{Ujemny rozkład a priori Dirichleta i eliminacja komponentów}

W~celu automatycznego określenia optymalnej liczby komponentów $M$ dla każdego piksela, Zivkovic
wprowadził do estymacji parametrów metodę maksymalizacji prawdopodobieństwa a posteriori (MAP)
z~użyciem sprzężonego rozkładu a priori Dirichleta~\cite{zivkovic2004, zivkovic2006}. Zdefiniowano
rozkład a priori w~postaci~\cite{zivkovic2004}:

\begin{equation}
  P = \prod_{m=1}^M \pi_m^{c_m}
  \label{eq:dirichlet_prior}
\end{equation}

Kluczowym elementem tej metody jest nadanie współczynnikom rozkładu a priori wartości ujemnych
($c_m = -c$)~\cite{zivkovic2004}. Z~matematycznego punktu widzenia ujemne współczynniki wprowadzają
karę za nadmierną złożoność modelu, co jest ściśle powiązane z~kryterium minimalnej długości opisu
(\textit{ang. Minimum Message Length} -- MML)~\cite{zivkovic2004}. Rozwiązanie równania MAP przy
uwzględnieniu ujemnego rozkładu a priori prowadzi do wyznaczenia wag
skorygowanych~\cite{zivkovic2004}:

\begin{equation}
  \hat{\pi}_m^{(t)} = \frac{1}{K} \left( \sum_{i=1}^t o_m^{(i)} - c \right)
  \label{eq:map_weights}
\end{equation}

gdzie parametr normalizacyjny wynosi $K = t - Mc$~\cite{zivkovic2004}. Aby zapobiec zanikaniu
wpływu kary wraz z~upływem czasu (gdy $t \to \infty$), ułamek $c/t$ zastępuje się stałą wartością
$c_T = c/T$~\cite{zivkovic2004}. Przy założeniu, że rzeczywista liczba komponentów opisujących
piksel jest niewielka, ostateczne równanie aktualizacji wag przyjmuje postać~\cite{zivkovic2004}:

\begin{equation}
  \hat{\pi}_m \leftarrow \hat{\pi}_m + \alpha (o_m^{(t)} - \hat{\pi}_m) - \alpha c_T
  \label{eq:mog2_final_weight_update}
\end{equation}

Zastosowanie stałej poprawki ujemnej $-\alpha c_T$ sprawia, że wagi komponentów, które nie są
regularnie wspierane przez nowe obserwacje, sukcesywnie maleją~\cite{zivkovic2004}. Gdy waga danego
komponentu osiągnie wartość ujemną, zostaje on trwale usunięty z~mieszaniny~\cite{zivkovic2004}.
W~ten sposób piksele reprezentujące stabilne elementy tła są opisywane jednym rozkładem Gaussa,
natomiast piksele w~obszarach o~złożonej dynamice (np. migoczące światła, falująca roślinność)
utrzymują większą liczbę komponentów~\cite{zivkovic2004, zivkovic2006}.

\subsection{KNN \textit{(ang. K-nearest Neighbors)}}

Algorytm KNN w kontekście odejmowania tła stanowi alternatywne, nieparametryczne podejście do
modelowania rozkładu prawdopodobieństwa pikseli. W podejściu nieparametrycznym funkcja gęstości
jest szacowana bezpośrednio na podstawie danych historycznych, bez zakładania konkretnego kształtu
rozkładu~\cite{elgammal2000}. Podstawową techniką wykorzystywaną w tym podejściu jest estymacja
gęstości jądrowej (KDE, \textit{ang. Kernel Density Estimation}), która wymaga przechowywania
zbioru $N$ próbek historycznych $\vec{x}_n$ reprezentujących wcześniej obserwowane wartości
pikseli~\cite{elgammal2000}.

\subsubsection{Estymacja gęstości jądrowej z użyciem estymatora balonowego}

Podstawowy estymator Parzena–Rosenblatta szacuje gęstość prawdopodobieństwa punktu $\vec{x}$ w
$d$-wymiarowej przestrzeni cech na podstawie $N$ próbek historycznych $\vec{x}_n$ jako:

\begin{equation}
  \hat{p}(\vec{x}) = \frac{1}{N h^d} \sum_{n=1}^{N} K\!\left(\frac{\vec{x} - \vec{x}_n}{h}\right)
\end{equation}

gdzie $K(\cdot)$ jest funkcją jądrową, a $h$ stałą szerokością pasma. W kontekście odejmowania tła
szerokość pasma decyduje o wrażliwości klasyfikacji: jeżeli piksel znajduje się poza promieniem
wyznaczonym przez jądro, zostaje zaklasyfikowany jako pierwszy plan~\cite{zivkovic2006}. Użycie
stałej szerokości pasma może generować błędy w obszarach o zróżnicowanej gęstości danych —
nadmierne wygładzenie w regionach gęstych i niedoszacowanie w rzadkich~\cite{terrell1992}.

Zastosowanie estymatora balonowego (\textit{ang. balloon estimator})~\cite{terrell1992} rozwiązuje
ten problem poprzez uzależnienie szerokości jądra $h$ od lokalnej gęstości w punkcie
estymacji~$\vec{x}$. Szerokość okna dobiera się odwrotnie proporcjonalnie do lokalnej gęstości:

\begin{equation}
  h(\vec{x}) = \frac{k}{\left[N\,\hat{p}(\vec{x})\right]^{1/d}}
\end{equation}

Wówczas estymowana gęstość prawdopodobieństwa wyraża się wzorem:

\begin{equation}
  \hat{p}(\vec{x}) \approx \frac{1}{N \cdot h(\vec{x})^d}
  \sum_{n=1}^{N} K\!\left(\frac{\vec{x} - \vec{x}_n}{h(\vec{x})}\right)
\end{equation}

Dla optymalizacji wydajności obliczeniowej w systemach czasu rzeczywistego jako funkcję jądrową
$K(u)$ stosuje się jądro jednostajne (prostokątne)~\cite{terrell1992}. Wówczas estymacja gęstości
sprowadza się do zliczenia próbek $\vec{x}_n$ leżących wewnątrz sfery o promieniu $h(\vec{x})$
wyznaczonym przez odległość do $k$-tego sąsiada~\cite{zivkovic2006}:

\begin{equation}
  \hat{p}(\vec{x}) \approx \frac{k}{N \cdot V_d \cdot h(\vec{x})^d}
\end{equation}

gdzie $V_d$ oznacza objętość jednostkowej kuli w przestrzeni $d$-wymiarowej. Takie podejście
drastycznie skraca czas obliczeń w porównaniu z jądrem gaussowskim, przy zachowaniu adaptacyjności
szerokości pasma~\cite{zivkovic2006, mittal2004}.

\subsubsection{Metoda k-najbliższych sąsiadów w przestrzeni barw RGB}

W praktycznej implementacji algorytmu KNN klasyfikacja nowego piksela $\vec{x}(t)$ opiera się na
następującym schemacie decyzyjnym~\cite{zivkovic2006}:

\begin{enumerate}
  \item Zakłada się, że piksel jest reprezentowany jako wektor w przestrzeni barw RGB: $\vec{x}(t) =
          [R(t),\, G(t),\, B(t)]^T$.

  \item Oblicza się odległości euklidesowe $d_n$ między $\vec{x}(t)$ a każdą z $N$ próbek zgromadzonych w
        buforze historycznym tego piksela $\mathcal{X} = \{\vec{x}_1, \dots,
          \vec{x}_N\}$~\cite{zivkovic2006, cucchiara2003}.

  \item Zlicza się liczbę próbek $k^*$, dla których $d_n$ jest mniejsza od progu $D_{thr}$ (w OpenCV
        parametr \texttt{dist2Threshold} przechowuje kwadrat tej odległości):
        \begin{equation}
          k^* = \sum_{n=1}^{N} \mathbf{1}\!\left[d\!\left(\vec{x}(t),\,\vec{x}_n\right) < D_{thr}\right]
        \end{equation}
        gdzie $\mathbf{1}[\cdot]$ oznacza funkcję wskaźnikową. Próbki spełniające ten warunek
        mieszczą się wewnątrz sfery o promieniu $D_{thr}$ wokół $\vec{x}(t)$~\cite{prati2003,
          zivkovic2006, opencv2025}.

  \item Jeśli $k^* \ge k$, piksel zostaje zaklasyfikowany jako tło~\cite{zivkovic2006, opencv2025}. W
        przeciwnym razie jest oznaczany jako pierwszy plan~\cite{zivkovic2006, opencv2025}.
\end{enumerate}

Aktualizacja modelu w algorytmie KNN zachodzi w sposób ciągły — nowa próbka piksela zastępuje
najstarszą próbkę w buforze (schemat FIFO). Dzięki temu model stopniowo dostosowuje się do wolnych
zmian warunków oświetleniowych, takich jak ruch chmur czy dynamika cieni~\cite{zivkovic2006} .

\subsection{Ocena efektywności algorytmów}

Wyniki segmentacji pierwszego planu weryfikuje się przez porównanie wygenerowanej maski binarnej
z~maską referencyjną (\textit{ang. ground truth}) na poziomie pikseli~\cite{cdnet2014}. Każdy
piksel klasyfikowany jest jako jeden z~czterech przypadków: prawdziwy pozytyw (TP -- piksel
pierwszego planu wykryty poprawnie), fałszywy pozytyw (FP -- piksel tła błędnie oznaczony jako
pierwszy plan), prawdziwy negatyw (TN -- piksel tła wykryty poprawnie) oraz fałszywy negatyw (FN --
piksel pierwszego planu pominięty przez algorytm). Na tej podstawie wyznacza się następujące
metryki~\cite{cdnet2014}:

\begin{itemize}
  \item \textbf{Czułość} \textit{(ang. Recall / True Positive Rate)} — stosunek
        poprawnie wykrytych pikseli pierwszego planu do wszystkich pikseli pierwszego
        planu w masce referencyjnej:
        \begin{equation}
          \mathrm{Recall} = \frac{TP}{TP + FN}
        \end{equation}

  \item \textbf{Swoistość} \textit{(ang. Specificity / True Negative Rate)} — stosunek
        poprawnie sklasyfikowanych pikseli tła do wszystkich pikseli tła w~masce
        referencyjnej:
        \begin{equation}
          \mathrm{Specificity} = \frac{TN}{TN + FP}
        \end{equation}

  \item \textbf{Precyzja} \textit{(ang. Precision)} — stosunek poprawnie wykrytych
        pikseli pierwszego planu do wszystkich pikseli oznaczonych przez algorytm jako
        pierwszy plan:
        \begin{equation}
          \mathrm{Precision} = \frac{TP}{TP + FP}
        \end{equation}

  \item \textbf{Miara F1} \textit{(ang. F1-score)} — średnia harmoniczna precyzji
        i~czułości, zapewniająca zrównoważoną ocenę przy nierównomiernym rozkładzie
        klas:
        \begin{equation}
          F_1 = 2 \cdot \frac{\mathrm{Precision} \cdot \mathrm{Recall}}
          {\mathrm{Precision} + \mathrm{Recall}}
        \end{equation}

  \item \textbf{IoU} \textit{(ang. Intersection over Union)} — stosunek pola
        przecięcia maski wykrytej i~referencyjnej do pola ich sumy. W~kontekście
        ewaluacji BGS operuje się na maskach binarnych, nie na ramkach
        otaczających~\cite{cdnet2014}:
        \begin{equation}
          \mathrm{IoU} = \frac{TP}{TP + FP + FN}
        \end{equation}

  \item \textbf{PWC} \textit{(ang. Percentage of Wrong Classifications)} — odsetek
        błędnie sklasyfikowanych pikseli względem całkowitej liczby pikseli w~sekwencji;
        stosowany jako metryka syntetyczna w~benchmarku CDnet~\cite{cdnet2014}:
        \begin{equation}
          \mathrm{PWC} = 100 \cdot \frac{FP + FN}{TP + FP + TN + FN}
        \end{equation}
\end{itemize}

Wymienione metryki są standardowo stosowane w~benchmarku CDnet 2014~\cite{cdnet2014}, który stanowi
punkt odniesienia dla porównania algorytmów BGS w~niniejszej pracy. Metryki oceny śledzenia wielu
obiektów (MOTA, HOTA, IDF1) oraz reidentyfikacji (Rank-1/CMC) zostają omówione w~kolejnych
podrozdziałach poświęconych odpowiednio modułom detekcji i~śledzenia.

\subsection{Analiza porównawcza i ewaluacja efektywności}

\subsubsection{Porównanie wydajności algorytmów BGS w trudnych warunkach atmosferycznych}

Efektywność algorytmów odejmowania tła zależy od charakterystyki badanej sceny oraz rodzaju
zakłóceń środowiskowych~\cite{bouwmans2014}. Tabela~\ref{tab:iou_bgs_comparison} przedstawia
średnie wartości współczynnika IoU uzyskane przez wybrane algorytmy BGS w~badaniach porównawczych
dla różnego rodzaju zakłóceń atmosferycznych wideo~\cite{karbowiak2022}.

\begin{table}[htbp]
  \centering
  \caption{Średnie wartości współczynnika IoU wybranych algorytmów BGS. Żródło: opracowanie własne na podstawie danych z~\cite{karbowiak2022}.}
  \label{tab:iou_bgs_comparison}
  \begin{tabular}{|l|c|c|c|c|c|c|}
    \hline
    \textbf{Środowisko / Warunki}
            & \textbf{CNT} & \textbf{KNN} & \textbf{MOG} & \textbf{MOG2} & \textbf{GMG} & \textbf{SOM} \\ \hline
    Śnieg   & 0,41         & 0,38         & 0,44         & 0,25          & 0,42         & 0,52         \\ \hline
    Wiatr   & 0,13         & 0,13         & 0,33         & 0,11          & 0,12         & 0,29         \\ \hline
    Deszcz  & 0,17         & 0,14         & 0,37         & 0,10          & 0,31         & 0,25         \\ \hline
    Średnia & 0,24         & 0,22         & 0,38         & 0,15          & 0,28         & 0,35         \\ \hline
  \end{tabular}
\end{table}

Analiza wyników przedstawionych w~Tabeli~\ref{tab:iou_bgs_comparison} pozwala na sformułowanie
kilku obserwacji dotyczących zachowania algorytmów BGS w~trudnych warunkach atmosferycznych.

Najsłabszą i~najbardziej równomierną odporność na zakłócenia środowiskowe wykazują algorytmy KNN
oraz MOG2 -- oba uzyskują najniższe wartości IoU zarówno przy wietrze (0,13 i~0,11), jak i~przy
deszczu (0,14 i~0,10). Algorytm MOG z~kolei wykazuje zaskakująco wysoką odporność w~warunkach
wietrznych (IoU~=~0,33) oraz deszczowych (IoU~=~0,37), co wynika z~jego większej bezwładności
modelu -- stała liczba komponentów gaussowskich wolniej absorbuje szybkie
zakłócenia~\cite{stauffer1999}.

Spośród wszystkich badanych metod najwyższą średnią skuteczność osiąga SOM (IoU~=~0,47),
a~najniższą -- MOG2 (IoU~=~0,25). Algorytm SOM (\textit{ang. Self-Organizing Map}) modeluje tło za
pomocą samoorganizacji poprzez szczutczną sieć neuronową~\cite{karbowiak2022}. Dzięki temu
skuteczniej radzi sobie ze złożonymi problemami, takimi jak ruszające się tło czy zmiany w
oświetleniu~\cite{karbowiak2022}.

Algorytm GMG osiąga wysoką średnią (IoU~=~0,42) dzięki bayesowskiemu modelowi tła budowanemu na
podstawie pierwszych klatek sekwencji (domyślnie 120 klatek w~OpenCV). Wysoka skuteczność przy
śniegu (IoU~=~0,42) jest efektem dobrego odwzorowania statystycznej zmienności tła w~fazie
inicjalizacji. Metoda ta ma jednak istotne ograniczenie -- wymaga stabilnej fazy uczenia i~słabo
radzi sobie z~dynamicznymi scenami od pierwszej klatki~\cite{karbowiak2022}.

Algorytm CNT (\textit{ang. Counter}, \textit{ang. Count of Non-stationary Targets}) opiera się na
zliczaniu, ile razy wartość piksela zmieniła się w danym czasie. Jest to metoda prosta
obliczeniowo, osiągająca zadowalające wyniki dla śniegu (IoU~=~0,41), jest zawodna przy wietrze
i~deszczu (IoU~=~0,13 i~0,17).

Pomimo że algorytmy SOM, GMG i~CNT uzyskują wyższe wartości IoU od KNN i~MOG2, w~niniejszej pracy
zdecydowano się na implementację tych ostatnich. Wynika to z~ich natywnej dostępności i~pełnego
wsparcia w~bibliotece OpenCV, niskiego zapotrzebowania na zasoby obliczeniowe niezbędnego
w~przetwarzaniu strumienia wideo w~czasie rzeczywistym, a~także szczegółowej dokumentacji
i~ugruntowanej pozycji w~literaturze~\cite{zivkovic2006}. Algorytmy SOM i~GMG wymagają odpowiednio
dłuższej fazy inicjalizacji lub wyższych zasobów, co utrudnia ich zastosowanie w~scenariuszach
z~dynamicznie zmieniającym się tłem bez etapu wstępnego uczenia.

\subsubsection{Wydajność obliczeniowa, ograniczenia sprzętowe i zjawisko dryfu tła}

W systemach czasu rzeczywistego wydajność obliczeniowa jest istotnym czynnikiem przy wdrażaniu
oprogramowania~\cite{bouwmans2014}. MOG2, dzięki dynamicznemu doborowi liczby komponentów
gaussowskich opartemu na kryterium MML \textit{(ang. Minimum Message Length)}~\cite{zivkovic2004},
charakteryzuje się niskim zużyciem pamięci RAM. Pozwala to na jego stosowanie w~architekturach
wielokamerowych na platformach wbudowanych~\cite{waghmare2026}.

Najpoważniejszą wadą klasycznych algorytmów BGS, w~tym MOG2 i~KNN, jest podatność na zjawisko dryfu
modelu tła (\textit{ang. background drift})~\cite{waghmare2026}. Jeżeli obiekt pierwszego plany
pozostanie nieruchomy, w ciągu 100 klatek może zostać zaklasyfikowany jako część
tła~\cite{waghmare2026}. Skutkuje to fragmentacją ścieżki śledzenia i~utratą ciągłości detekcji.
Dlatego na trudnych zbiorach testowych samodzielne algorytmy BGS osiągają niską czułość
detekcji~\cite{waghmare2026}.

%%%%%%%%%%%%%%%%%%%%%%%%%%%%%%%%%%%%%%%%%%%%%%%%%%%%%%%%%%%
%%%%%%%%%%%%%%%%%%%%%%%%%%%%%%%%%%%%%%%%%%%%%%%%%%%%%%%%%%%

\section{Detekcja obiektów oparta na głębokim uczeniu (Deep Learning)}

Przed upowszechnieniem wielowarstwowych sieci neuronowych, systemy detekcji obiektów opierały się
na ręcznie projektowanych cechach wizualnych \textit{(ang. hand-crafted features)}, które były
następnie klasyfikowane za pomocą tradycyjnych algorytmów uczenia maszynowego. Jednym
z~najważniejszych rozwiązań w~obszarze detekcji obiektów w~czasie rzeczywistym stał się algorytm
Viola-Jones, pierwotnie opracowany do detekcji twarzy. Metoda ta opiera się na trzech kluczowych
koncepcjach: reprezentacji obrazu całkowego \textit{(ang. Integral Image)}, wyznaczaniu prostych
cech \textit{(ang. Haar-like features)} oraz kaskadowej strukturze klasyfikatorów opartej na
algorytmie AdaBoost~\cite{viola2001}. Zastosowanie obrazu całkowego umożliwiło obliczanie sum
jasności pikseli wewnątrz prostokątnych obszarów w czasie stałym $O(1)$, niezależnie od skali.
Pomimo rewolucyjnej szybkości działania, detektor Viola-Jones wykazywał niską odporność na rotacje
obiektów w przestrzeni trójwymiarowej, okluzje oraz gwałtowne zmiany kontrastu i oświetlenia
sceny~\cite{viola2001}. Kolejny przełomowy algorytm, służący do detekcji sylwetek ludzkich, to
metoda HOG~+~SVM. Algorytm opierał się na histogramach zorientowanych gradientów liczonych w
komórkach obrazu i normalizowanych lokalnie w nakładających się blokach, co poprawiało odporność na
zmiany kontrastu~\cite{dalal2005}. Autorzy wykazali, że taka reprezentacja znacznie przewyższa
wcześniejsze zestawy cech w zadaniu detekcji ludzi. Ograniczeniem podejścia pozostawała m.in.
podatność na deformacje postawy oraz koszt obliczeniowy wynikający ze skanowania obrazu metodą
przesuwanego okna~\cite{dalal2005}.

\subsection{Splotowe sieci neuronowe (CNN) i detektory dwuetapowe}
Rok 2012 przyniósł zasadnicze zmiany w wizji komputerowej wraz z sukcesem głębokiej sieci splotowej
AlexNet w konkursie ILSVRC~\cite{krizhevsky2012imagenet}. Model ten, zbudowany z pięciu warstw
splotowych i trzech warstw w pełni połączonych, pokazał, że sieci neuronowe mogą samodzielnie uczyć
się użytecznych cech obrazu lepiej niż metody ręcznie projektowane~\cite{krizhevsky2012imagenet}.
Zastosowanie funkcji ReLU oraz technik regularizacji, takich jak dropout, umożliwiło skuteczne
trenowanie dużych modeli na procesorach graficznych (GPU).

Zastosowanie sieci splotowych do lokalizacji obiektów doprowadziło do powstania detektorów
dwuetapowych. W pierwszym kroku system wskazywał możliwe miejsca występowania obiektu, a w drugim
klasyfikował je i doprecyzowywał położenie. Pierwszym istotnym modelem był R-CNN, który łączył
tradycyjne generowanie propozycji regionów, na przykład metodą Selective Search, z siecią CNN
analizującą każdy region osobno~\cite{girshick2014rich}. Było to jednak wolne, ponieważ ten sam
fragment obrazu był przetwarzany wiele razy. Fast R-CNN przyspieszył ten proces, obliczając cechy
dla całego obrazu tylko raz i wykorzystując warstwę ROI Pooling do pobierania informacji dla
proponowanych regionów~\cite{girshick2015fast}. Dalszym ulepszeniem był Faster R-CNN, który
zastąpił Selective Search siecią RPN, czyli modułem proponującym regiony bezpośrednio na podstawie
map cech~\cite{ren2015faster}. Dzięki temu model stał się znacznie szybszy i bardzo dokładny, choć
nie zawsze osiągał pełny czas rzeczywisty.

\subsection{Detektory jednoetapowe}

Modele dwuetapowe, mimo wysokiej dokładności, często były zbyt wolne do zastosowań wymagających
szybkiej reakcji, takich jak monitoring czasu rzeczywistego czy systemy wbudowane. Odpowiadając na
tę potrzebę,powstały detektory jednoetapowe (ang. \textit{one-stage detectors}), które bezpośrednio
przekształcają obraz wejściowy w predykcje klas oraz współrzędnych ramek otaczających, bez osobnego
etapu generowania propozycji regionów.

Jednym z najważniejszych algorytmów tej grupy stał się YOLO \textit{(ang. You Only Look Once)}, w
którym potraktowano detekcję obiektów jako pojedynczy problem regresji realizowany przez jedną sieć
neuronową w jednym przebiegu obrazu (\textit{single forward pass})~\cite{redmon2016yolo}. Dzięki
temu możliwe stało się wykrywanie obiektów w czasie zbliżonym do rzeczywistego, przy zachowaniu
dobrej skuteczności. Kolejne wersje tej rodziny rozwijały to podejście, poprawiając kompromis
między szybkością a dokładnością. YOLOv7 wprowadził architekturę E-ELAN \textit{(ang. Extended
  Efficient Layer Aggregation Network)}, która usprawniła przepływ informacji w sieci, oraz koncepcję
\textit{compound scaling}, pozwalającą lepiej dobierać rozmiar modelu do możliwości obliczeniowych.
Autorzy zastosowali także tzw. \textit{bag-of-freebies}, czyli zestaw ulepszeń poprawiających
jakość uczenia bez zwiększania kosztu inferencji~\cite{wang2023yolov7}. YOLOv8 przeszedł na
podejście bezkotwicowe (\textit{anchor-free}), rezygnując z wcześniej definiowanych ramek bazowych,
oraz zastosował rozdzieloną głowę detekcyjną \textit{(ang. split head)}, w której klasyfikacja i
regresja są wykonywane osobno~\cite{ultralytics2023yolov8}. Z kolei YOLOv9 zaproponował rozwiązania
mające ograniczać utratę informacji w głębokich sieciach oraz poprawiać przepływ gradientu, co ma
szczególne znaczenie w przypadku małych i trudnych do wykrycia obiektów.~\cite{wang2024yolov9}.

Alternatywnym podejściem do detekcji jednoetapowej jest SSD \textit{(ang. Single Shot MultiBox
  Detector)}. Metoda ta wykorzystuje kilka map cech o różnych rozdzielczościach, dzięki czemu lepiej
radzi sobie z obiektami o zróżnicowanych rozmiarach~\cite{liu2016ssd}. W porównaniu z modelami
dwuetapowymi SSD oferował prostszą architekturę i większą szybkość działania, choć początkowo
często ustępował im dokładnością.

Głównym problemem detektorów jednoetapowych była jednak silna nierównowaga klasowa: w obrazie
dominuje tło, a obiekty zajmują tylko niewielką część danych. Prowadziło to do sytuacji, w której
model był zbyt mocno uczony rozpoznawania tła, a zbyt słabo skupiał się na rzeczywistych
obiektach~\cite{lin2017focal}. Problem ten skutecznie złagodził RetinaNet, który wprowadził funkcję
straty Focal Loss:
\[
  FL(p_t)=-\alpha_t(1-p_t)^\gamma\log(p_t)
\]
Czynnik modulujący $(1 - p_t)^\gamma$ (gdzie $\gamma$ to parametr skupienia) sprawia, że w tym
ujęciu łatwe przykłady, czyli takie, które model klasyfikuje z dużą pewnością, otrzymują mniejszą
wagę w procesie uczenia. Dzięki temu gradient jest bardziej skoncentrowany na trudnych przykładach
i na obiektach pierwszoplanowych, co poprawia skuteczność detekcji bez rezygnacji z szybkości
charakterystycznej dla podejścia jednoetapowego~\cite{lin2017focal}.

\subsection{Transformery w detekcji obiektów}

Wprowadzenie mechanizmów uwagi, znanych wcześniej głównie z przetwarzania języka naturalnego,
doprowadziło do powstania architektury DETR (\textit{Detection Transformer}). Model ten traktuje
detekcję obiektów jako problem bezpośredniego przewidywania zbioru wyników (\textit{direct set
  prediction}). W praktyce oznacza to, że zamiast stosować kotwice (\textit{anchors}) i dodatkowy
etap usuwania nadmiarowych ramek NMS \textit{(ang. Non-Maximum Suppression)}, DETR łączy ekstraktor
cech oparty na sieci splotowej z enkoderem i dekoderem Transformer oraz mechanizmem dopasowania
dwudzielnego opartym na algorytmie węgierskim~\cite{carion2020end}.

Mimo eleganckiej konstrukcji oryginalny DETR miał istotne ograniczenia. Model uczył się bardzo
wolno, wymagał dużych zasobów obliczeniowych i osiągał słabsze wyniki w detekcji małych obiektów.
Odpowiedzią na te problemy był Deformable DETR. Zastosowano w nim odkształcalną uwagę
(\textit{deformable attention}), która zamiast analizować cały obraz, skupia się na niewielkiej
liczbie punktów wybranych wokół punktów referencyjnych. Dzięki temu model lepiej wykorzystuje cechy
wieloskalowe, szybciej się uczy i radzi sobie skuteczniej z obiektami o różnych rozmiarach
~\cite{zhu2021deformable}.

Nowszym rozwinięciem tej linii jest RT-DETR (\textit{Real-Time Detection Transformer}),
zaprojektowany z myślą o zastosowaniach czasu rzeczywistego. Model ten łączy wydajny enkoder
hybrydowy z mechanizmem wyboru najbardziej obiecujących zapytań dekodera, co pozwala ograniczyć
koszty obliczeniowe i zachować wysoką szybkość działania. W odróżnieniu od klasycznych detektorów
opartych na CNN RT-DETR dąży do osiągnięcia bardzo dobrego kompromisu między dokładnością a czasem
inferencji, a przy tym eliminuje potrzebę stosowania NMS jako osobnego etapu przetwarzania
~\cite{zhao2024detrs}.

\subsection{Dane treningowe a problem luki domenowej}

Skuteczność głębokich modeli detekcji zależy od wielkości i jakości danych uczących. Modele są
zazwyczaj trenowane na dużych, ogólnych zbiorach badawczych, takich jak MS COCO lub ImageNet.
Jednak bezpośrednie przeniesienie modelu wytrenowanego na danych syntetycznych do rzeczywistego
środowiska przemysłowego lub systemu CCTV zawodzi, ujawniając problem luki domenowej \textit{(ang.
  domain shift)}\~\cite{he2024}.

Luka domenowa wynika z różnic w charakterystyce obrazów między zbiorem treningowym a rzeczywistym
środowiskiem.
\begin{itemize}
  \item Zdjęcia w zbiorach akademickich najczęściej wykonywane są z perspektywy na wysokości oczu, kamery
        monitoringu są zwykle umieszczone znacznie wyżej. Powoduje to zniekształcenia perspektywy, które
        wpływają na kształt i rozmiar obiektów, a także na ich wzajemne relacje przestrzenne.
  \item Kamery CCTV charakteryzują się specyficznymi zniekształceniami optycznymi (np. obiektywy
        szerokokątne typu rybie oko), silną kompresją stratną strumienia wideo oraz wysokim poziomem szumów
        matrycy w warunkach nocnych.
  \item Zbiory testowe rzadko reprezentują skrajne zjawiska pogodowe (silny deszcz, śnieżyca, mgła,
        oślepiające słońce generujące głębokie cienie), które w rzeczywistym monitoringu zewnętrznym
        występują powszechnie~\cite{he2024}.
\end{itemize}

Najczęściej stosowane techniki niwelowania luki domenowej w detekcji obiektów można podzielić na
trzy główne grupy: translację domen, samouczenie z pseudoetykietami oraz dopasowanie cech między
domenami. Pierwsza grupa polega na przekształcaniu danych źródłowych tak, aby bardziej przypominały
dane docelowe, zwykle przy użyciu modeli generatywnych. Wiąże się to jednak z utratą informacji i
wysoką złożonością obliczeniową.~\cite{he2024}.

Metody samouczenia \textit{(ang. self-training)}, w których model uczony jest na danych źródłowych
z etykietami oraz na danych docelowych opisanych pseudoetykietami generowanymi przez model
nauczyciela. Jego skuteczność zależy jednak od jakości pseudoetykiet, dlatego wiele nowszych metod
wprowadza mechanizmy ich selekcji, ważenia lub filtrowania~\cite{he2024}.

Szeroko stosowane metody dopasowania cech (\textit{domain alignment}), dążą do zmniejszenia różnic
pomiędzy rozkładami cech w domenie źródłowej i docelowej. W klasycznych rozwiązaniach wykorzystuje
się do tego sieci przeciwstawne, w których dyskryminator próbuje odróżnić cechy pochodzące z obu
domen, a ekstraktor cech uczy się tworzyć reprezentacje nieodróżnialne z punktu widzenia domeny. W
detekcji obiektów takie dopasowanie może zachodzić na poziomie obrazu, regionu lub pojedynczej
instancji~\cite{chen2018, he2024}.

Nowsze badania wskazują jednak, że globalne wyrównywanie cech na poziomie całego obrazu szkodzi
precyzji lokalizacji detektorów, gdyż prowadzi do transferu informacji z nieistotnych obszarów tła
~\cite{he2024}. W związku z tym wprowadzono zaawansowane strategie adaptacji różnicowej (np. UFOA –
Uncertainty-based Foreground-Oriented Alignment)~\cite{he2024}, które selektywnie nadają wyższy
priorytet obszarom reprezentującym obiekty pierwszoplanowe.

\subsection{Detekcja w trudnych warunkach}

Wdrożenie skutecznych algorytmów detekcji obiektów w systemach monitoringu wymaga wysokiej
odporności modeli na zakłócenia. Do najczęstszych problemów należą okluzje, gwałtowne zmiany
oświetlenia oraz niska rozdzielczość małych lub odległych obiektów~\cite{nikouei2025, cheng2026}.

Techniki ograniczające wpływ okluzji obejmują między innymi:
\begin{itemize}
  \item \textbf{Generatywne warstwy kompozycyjne} – w tym podejściu klasyczne głowice klasyfikacyjne zastępowane są strukturą generatywną, która potrafi wnioskować o niewidocznych częściach obiektu na podstawie relacji przestrzennych między jego widocznymi fragmentami. Przykładem takiego rozwiązania jest model CompositionalNet \cite{cheng2026}.
  \item \textbf{Algorytmy głosowania częściowego \textit{(ang. part-based voting)}} – model uczy się niezależnych reprezentacji poszczególnych części obiektu, a następnie szacuje położenie całego obiektu na podstawie wykrytych elementów \cite{cheng2026}. W detekcji pieszych może to oznaczać osobne wykrywanie głowy, ramion i nóg, a następnie wyznaczenie położenia całej sylwetki.
  \item \textbf{Łączenie wielu widoków i wspólne wnioskowanie} – w systemach wielokamerowych obiekt zasłonięty z jednej perspektywy może być dobrze widoczny z innej. Połączenie predykcji z kilku sąsiadujących kamer pozwala trafniej odtworzyć położenie obiektu, nawet jeśli w pojedynczym obrazie jest on silnie zasłonięty \cite{cheng2026}.
\end{itemize}

Systemy CCTV działają przez całą dobę, dlatego są narażone na zmienne warunki oświetleniowe
utrudniające detekcję. Problem stanowi zarówno intensywne światło słoneczne powodujące głębokie
cienie, jak i nagłe prześwietlenia, na przykład od reflektorów samochodowych w nocy, a także bardzo
słabe oświetlenie po zmroku~\cite{nikouei2025}. Dodatkowo opady atmosferyczne, takie jak deszcz i
śnieg, wprowadzają zakłócenia do obrazu i osłabiają składowe wysokoczęstotliwościowe, co prowadzi
do rozmycia krawędzi obiektów~\cite{wille2026}. Jednym ze sposobów ograniczania utraty detali są
moduły o zmiennym polu recepcyjnym, takie jak RepHELAN, które wykorzystują heterogeniczne
konwolucje do zwiększania efektywnego pola recepcyjnego bez pogorszenia reprezentacji małych
obiektów~\cite{yang2024}. Z kolei w trudnych warunkach oświetleniowych stosuje się mechanizmy uwagi
wieloskalowej, na przykład LSKA (\textit{Large Separable Kernel Attention}), które mogą
adaptacyjnie ograniczać wpływ lokalnych prześwietleń i koncentrować się na bardziej stabilnych
zależnościach przestrzennych~\cite{tang2026}.

W rozległych scenach monitoringu, takich jak place miejskie, skrzyżowania czy parkingi, obiekty
znajdujące się daleko od kamery klasyfikowane są jako małe jeżeli zajmują na obraziemniej niż $32
  \times 32$, a~często stanowią nawet mniej niż $20 \times 20$ pikseli~\cite{nikouei2025}.
W~klasycznych sieciach konwolucyjnych kolejne operacje zmniejszania rozdzielczości, realizowane na
przykład przez konwolucje ze skokiem lub warstwy MaxPooling, mogą prowadzić do utraty informacji o
tak małych obiektach jeszcze przed przekazaniem cech do głowic predykcyjnych~\cite{nikouei2025}.
Aby temu przeciwdziałać, nowoczesne architektury wykorzystują mechanizmy wzmacniające sygnał małych
obiektów. Jednym z nich są gęste połączenia wspomagające, takie jak moduł SAF (\textit{Superficial
  Assisted Fusion}) stosowany w części agregującej cechy (\textit{neck}), który przekazuje do
głębszych warstw informacje o wysokiej rozdzielczości przestrzennej z początkowych etapów
ekstrakcji cech~\cite{yang2024}. Innym rozwiązaniem jest regresja dystrybucyjna, na przykład DFL
(\textit{Distribution Focal Loss}) wykorzystywana w modelach YOLOv8 i YOLOv9, w której współrzędne
ramki nie są modelowane jako pojedyncze wartości, lecz jako rozkłady prawdopodobieństwa wokół jej
krawędzi. Pozwala to precyzyjniej opisać niepewność lokalizacji małego obiektu, także na poziomie
subpikselowym~\cite{yang2024}.
%%%%%%%%%%%%%%%%%%%%%%%%%%%%%%%%%%%%%%%%%%%%%%%%%%%%%%%%%%%
%%%%%%%%%%%%%%%%%%%%%%%%%%%%%%%%%%%%%%%%%%%%%%%%%%%%%%%%%%%

\section{Śledzenie wielu obiektów w pojedynczej kamerze}

Śledzenie wielu obiektów (Multiple Object Tracking – MOT) w strumieniu wideo z pojedynczej kamery polega na jednoczesnym szacowaniu pozycji wielu poruszających się obiektów w czasie oraz na zachowaniu ich spójnej tożsamości w kolejnych klatkach sekwencji obrazów. Zadanie to wymaga rozwiązania trzech podstawowych problemów:
\begin{itemize}
  \item \textbf{Asocjacji danych} \textit{(ang. data association)} - powiązania niezależnych detekcji tych samych obiektów na przestrzeni kolejnych klatek wideo i przypisania im stałego, unikalnego identyfikatora ang. \textit{(ang. track ID)} \cite{bewley2016}.
  \item \textbf{Inicjalizacji nowych ścieżek} — wykrywania momentu wejścia nowego, dotąd nieśledzonego obiektu w pole widzenia kamery i utworzenia dla niego nowej trajektorii \cite{bewley2016}.
  \item \textbf{Terminacji ścieżek} — identyfikacji momentu, w którym obiekt opuszcza monitorowaną scenę (lub zostaje długotrwale zasłonięty), w celu zakończenia aktualizacji jego trajektorii. Operacja ta zapobiega nieograniczonemu wzrostowi liczby śledzonych obiektów oraz redukuje błędy lokalizacji wynikające z długotrwałej predykcji ruchu \cite{bewley2016}.
\end{itemize}
\subsection{Paradygmat śledzenia ścieżek \textit{(ang. tracking-by-detection)}}
Od około 2015 roku w dziedzinie wieloobiektowego śledzenia dominuje paradygmat śledzenia przez detekcję \textit{(ang. tracking-by-detection lub detection-based tracking)} – DBT \cite{wojke2017}. Podejście to opiera się na rozdzieleniu modułu lokalizacji przestrzennej od modułu odpowiedzialnego za utrzymanie tożsamości. Przepływ informacji realizowany jest kaskadowo: najpierw algorytm detekcyjny (np. głęboka sieć neuronowa z rodziny YOLO) niezależnie lokalizuje obiekty w każdej klatce wideo, a następnie moduł śledzący łączy te dyskretne predykcje w ciągłe trajektorie \cite{bewley2016, wojke2017}. Taka separacja w architekturze systemu sprawia, że precyzja i czułość całego systemu MOT są bezpośrednio ograniczone przez jakość działania zastosowanego detektora bazowego \cite{bewley2016}.

\subsection{Metody predykcji stanu i algorytmy skojarzenia}
Aby efektywnie połączyć detekcje z klatki $t$ z obiektami w klatce $t+1$, nowoczesne systemy
śledzenia wykorzystują kombinację filtrów estymacji stanu oraz algorytmów optymalizacji dyskretnej.

\subsubsection{Filtr Kalmana}

Do modelowania dynamiki ruchu obiektów oraz predykcji ich trajektorii w~przestrzeni obrazu stosuje
się liniowy, dyskretny filtr Kalmana~\cite{kalman1960, bewley2016}. Zakładając model ruchu o~stałej
prędkości (\textit{ang. Constant Velocity Model}), filtr ten pozwala na optymalne szacowanie stanu
poruszającego się obiektu dla zaszumionych detekcji brzegowych. Zapewnia to ciągłość śledzenia
nawet podczas przejściowych zakłóceń lub krótkotrwałych okluzji~\cite{bewley2016}.

Matematyczny zapis opiera się na założeniu, że stan obiektu w~chwili $k$ określa równanie stanu,
a~jego sensoryczny odczyt definiuje równanie obserwacji~\cite{kalman1960}:
\begin{equation}
  x_k = F_k x_{k-1} + w_{k-1}
  \label{eq:state_equation}
\end{equation}
\begin{equation}
  z_k = H_k x_k + v_k
  \label{eq:measurement_equation}
\end{equation}
gdzie $x_k$ reprezentuje poszukiwany wektor stanu systemu, $z_k$ oznacza dostarczany przez detektor rzeczywisty wektor pomiarowy, natomiast $F_k$ i~$H_k$ to kolejno: macierz przejścia stanu oraz macierz obserwacji~\cite{kalman1960}.

Nieujemne zmienne losowe $w_k$ i~$v_k$ reprezentują niezależne od siebie szumy procesu oraz pomiaru
o~wielowymiarowym rozkładzie Gaussa:
\begin{equation}
  w_k \sim \mathcal{N}(0, Q_k), \quad v_k \sim \mathcal{N}(0, R_k)
\end{equation}
gdzie $Q_k$ jest macierzą kowariancji szumu procesu, odzwierciedlającą niedoskonałości przyjętego modelu dynamiki, a~$R_k$ oznacza macierz kowariancji szumu pomiarowego, wynikającą z~dokładności detektora~\cite{kalman1960}.

Algorytm filtru działa rekurencyjnie w dwuetapowej pętli sprzężenia zwrotnego, na którą składają
się faza predykcji (\textit{ang. predict}) oraz korekty (\textit{ang. update})~\cite{kalman1960,
  ryspekov2026}:

\begin{itemize}
  \item \textbf{Etap Predykcji} Na podstawie stanu oszacowanego w poprzednim kroku czasowym, model matematyczny rzutuje pozycję oraz kowariancję błędu obiektu w przyszłość. W tym kroku wyznaczana jest wstępna hipoteza dotycząca nowego położenia ramki pieszego, zanim jeszcze system przetworzy nową klatkę wideo. Pozwala to na zachowanie ciągłości trajektorii nawet w przypadku chwilowego braku odczytu z detektora.
  \item \textbf{Etap Korekty} Po zarejestrowaniu rzeczywistego pomiaru z detektora, wstępna predykcja zostaje zaktualizowana. Obliczana jest różnica między stanem przewidywanym a rzeczywistym, a następnie – przy użyciu optymalnego wzmocnienia Kalmana – następuje korekta wektora stanu i minimalizacja kowariancji błędu, co przygotowuje filtr do kolejnej iteracji.
\end{itemize}

\subsubsection{Algorytm węgierski}

Algorytm węgierski to globalny optymalizator kombinatoryczny. Jego celem jest rozwiązanie problemu
optymalnego skojarzenia (\textit{ang. data association}) wyznaczonych hipotez ruchu z nowo
zarejestrowanymi przez detektor pomiarami. Proces śledzenia zostaje sprowadzony do problemu
przepływu w grafie dwudzielnym~\cite{bewley2016, kuhn1955}.

Problem skojarzenia danych w klatce $t$ formułuje się jako zadanie liniowego przypisania
(\textit{ang. Linear Assignment Problem -- LAP}) o minimalnym koszcie~\cite{kuhn1955}. W tym celu
konstruuje się macierz kosztów $\mathbf{C}$ o wymiarach $N \times M$~\cite{ristani2016}:
\begin{equation}
  \mathbf{C} =
  \begin{bmatrix}
    c_{1,1} & c_{1,2} & \dots  & c_{1,M} \\
    c_{2,1} & c_{2,2} & \dots  & c_{2,M} \\
    \vdots  & \vdots  & \ddots & \vdots  \\
    c_{N,1} & c_{N,2} & \dots  & c_{N,M}
  \end{bmatrix}
\end{equation}
gdzie $N$ reprezentuje liczbę aktywnych trajektorii - przewidywanych przez filtr Kalmana, $M$ to liczba nowych detekcji w bieżącej klatce, a pojedynczy element $c_{i,j}$ określa koszt przypisania detekcji $j$ do trajektorii $i$~\cite{ristani2016}. W najprostszym wariancie, na przykład w algorytmie SORT~\cite{bewley2016}, koszt ten definiuje się jako odległość IoU:
\begin{equation}
  c_{i,j} = 1 - IoU(T_i, D_j)
\end{equation}
gdzie $T_i$ to przewidziany obszar trajektorii, a $D_j$ to obszar nowej detekcji~\cite{bewley2016}. W bardziej zaawansowanych systemach, takich jak BoT-SORT~\cite{aharon2022}, koszt ten łączony jest z odległością cosinusową wyodrbębnionych cech wizualnych Re-ID.

Cel optymalizacyjny algorytmu węgierskiego polega na znalezieniu takiej binarnej macierzy przypisań
$\mathbf{X}$ która minimalizuje sumaryczny koszt dopasowania i dla której $x_{i,j} = 1$, gdy
detekcja $j$ została przypisana do trajektorii $i$, oraz $x_{i,j} = 0$ w przeciwnym
razie~\cite{kuhn1955, ristani2016}:
\begin{equation}
  \min \sum_{i=1}^{N} \sum_{j=1}^{M} c_{i,j} x_{i,j}
\end{equation}
przy zachowaniu mapowania jeden do jednego, co oznacza, że jedna trajektoria może zostać przypisana do maksymalnie jednej nowej detekcji, a każda detekcja może obsłużyć co najwyżej jedną trajektorię~\cite{kuhn1955}

Dzięki temu sformułowaniu matematycznemu, integracja predykcji ruchu z filtru Kalmana oraz
globalnej optymalizacji algorytmu węgierskiego pozwala systemowi na precyzyjne odróżnianie blisko
położonych obiektów i mitygację błędów typu \textit{ID switch}.

\subsection{Ewolucja kluczowych algorytmów śledzenia}
Rozwój algorytmów klasy MOT doprowadził do powstania kilku przełomowych architektur, stopniowo
rozwiązujących problemy okluzji i wydajności obliczeniowej.

\subsubsection{SORT \textit{(ang. Simple Online and Realtime Tracking)}, 2016:}
Minimalistyczny algorytm łączący predykcję filtrem Kalmana i dopasowanie oparte na metryce IoU za
pomocą algorytmu węgierskiego~\cite{wojke2017}. Dzięki prostej budowie cechuje się dużą szybkością - aktualizuje
ścieżki z częstotliwościom 260 $Hz$~\cite{bewley2016}, jednak jest podatny na gubienie tożsamości w
przypadku nawet krótkotrwałego zasłonięcia obiektu~\cite{bewley2016, wojke2017}.

\subsubsection{Deep SORT, 2017}
Rozwinięcie algorytmu SORT o mechanizm asocjacji oparty na głębokich cechach
wizualnych~\cite{wojke2017}. Poprzez zastosowanie dedykowanej sieci CNN, wytrenowanej na zadanie
reidentyfikacji, Deep SORT wyznacza deskryptory wizualne obiektów. Pozwala to na poprawne
powiązanie tożsamości po długotrwałych okluzjach, drastycznie redukując liczbę błędnych przełączeń
tożsamości (ang. \textit{ID switches}) – w badaniach wykazano spadek tego błędu o 45\% w stosunku
do SORT~\cite{wojke2017}.

\subsubsection{ByteTrack, 2022}
Autorzy tego algorytmu~\cite{zhang2022} zwrócili uwagę na wadę wcześniejszych metod, polegającą na
odrzucaniu detekcji o niskim poziomie ufności - poniżej danego progu. Wskazali oni, że takie
pomijane detekcje w rzeczywistości bardzo często reprezentują prawdziwe obiekty, które uległy
częściowemu zasłonięciu lub rozmyciu w ruchu. Aby rozwiązać ten problem, wprowadzono dwustopniowy
mechanizm asocjacji BYTE. W pierwszym kroku system standardowo dopasowuje detekcje o wysokiej
ufności. Jeśli po tym etapie istnieją nieprzypisane trajektorie, algorytm przechodzi do drugiego
kroku. Zamiast usuwać słabe detekcje jako szum tła, system sprawdza, czy ich geometria pokrywa się
z miejscem, w którym filtr Kalmana przewidywał obecność śledzonego wcześniej obiektu. W tej fazie
celowo wykorzystuje się wyłącznie podobieństwo przestrzenne (metrykę IoU), ignorując cechy
wizualne, które przy silnej okluzji są niewiarygodne. Zaufanie słabej detekcji, o ile znajduje się
ona w przewidywanym miejscu, pozwala z sukcesem kontynuować śledzenie trudnych do wykrycia
obiektów, co ogranicza fragmentację ścieżek i redukuje liczbę błędów typu False Negative.

\subsubsection{BoT-SORT, 2022}
Zestaw optymalizacji (\textit{ang. bag of tricks}) zintegrowany z algorytmem ByteTrack, który
wprowadza trzy zmiany rozwiązujące główne ograniczenia wcześniejszych metod~\cite{aharon2022}:
\begin{itemize}
  \item \textbf{Kompensację ruchu kamery (CMC):} Wykorzystuje globalną kompensację ruchu (GMC) oraz rzadki przepływ optyczny do estymacji transformacji tła. Pozwala to na matematyczną korektę wektora stanu filtru Kalmana, zapobiegając rozbieżnościom i utracie ścieżek przy drganiach lub gwałtownych ruchach kamery.
  \item \textbf{Zmodyfikowany wektor stanu:} Wprowadza bezpośrednie modelowanie szerokości i wysokości ramki ograniczającej w filtrze Kalmana, rezygnując z powszechnie stosowanej miary proporcji boków. Zmiana ta poprawiła precyzję dopasowania predykcji do rzeczywistej szerokości obiektów.
  \item \textbf{Nowy mechanizm łączenia macierzy kosztów:} Integruje odległość przestrzenną IoU oraz unikalne cechy wyglądu Re-ID. Zamiast stosowanej powszechnie sumy ważonej, algorytm jako ostateczny koszt przypisania wybiera wartość minimalną z obu macierzy (po uprzednim odfiltrowaniu niemożliwych par), co zwiększa skuteczność kojarzenia obiektów po okluzjach.
\end{itemize}

\subsection{Ocena efektywności algorytmów MOT}

Ilościowa ocena systemów śledzenia wieloobiektowego (MOT) wymaga zastosowania ujednoliconych
wskaźników matematycznych, które weryfikują zarówno zdolność detekcji obiektów w poszczególnych
klatkach, jak i stabilność utrzymania ich tożsamości w czasie. Punktem wyjścia do wyznaczenia
zaawansowanych metryk jest analiza podstawowych błędów klasyfikacji, wyznaczanych na poziomie
pojedynczych klatek wideo po etapie asocjacji:

\begin{itemize}
  \item \textbf{False Negatives ($FN$):} Błędy pominięcia, określające sytuację, w której rzeczywisty obiekt obecny w kadrze nie został wykryty przez detektor lub nie przypisano go do żadnej aktywnej trajektorii.
  \item \textbf{False Positives ($FP$):} Błędy fałszywego alarmu, występują gdy system tworzy lub utrzymuje trajektorię w miejscu, w którym w rzeczywistości nie ma żadnego obiektu.
  \item \textbf{Identity Switches ($IDSW$):} Błędy zmiany tożsamości, opisują zdarzenie, w którym poprawnie śledzony obiekt otrzymuje od nowy, nieużywany dotąd identyfikator lub przejmuje identyfikator innego obiektu.
\end{itemize}

Na podstawie powyższych błędów oraz liczby poprawnie zidentyfikowanych obiektów (True Positives --
$TP$), wyznacza się elementarne wskaźniki skuteczności detektora: czułość systemu (\textit{ang.
  Recall}) oraz jego precyzję (\textit{ang. Precision}):
\begin{equation}
  Recall = \frac{TP}{TP + FN}, \quad Precision = \frac{TP}{TP + FP}
\end{equation}

Historycznie standardem w ewaluacji systemów śledzenia, zdefiniowanym w ramach metodologii CLEAR
MOT, są wskaźniki \textbf{MOTA} oraz \textbf{MOTP}~\cite{bernardin2008}:

\begin{itemize}
  \item \textbf{MOTA (Multiple Object Tracking Accuracy):} Ocenia ogólną dokładność systemu śledzącego, integrując trzy opisane wcześniej typy błędów w odniesieniu do sumarycznej liczby obiektów rzeczywistych (Ground Truth, $GT$) we wszystkich klatkach potoku wideo~\cite{bernardin2008, milan2016}:
        \begin{equation}
          MOTA = 1 - \frac{\sum_{t}(FN_{t} + FP_{t} + IDSW_{t})}{\sum_{t}GT_{t}}
        \end{equation}
        gdzie $t$ oznacza indeks klatki obrazu. Metryka ta przyjmuje wartości z przedziału $(-\infty, 1]$. Ze względu na fakt, że błędy detekcji ($FN$ i $FP$) są zazwyczaj liczniejsze niż błędy asocjacji ($IDSW$), wskaźnik MOTA faworyzuje jakość samego detektora kosztem stabilności trajektorii.

  \item \textbf{MOTP (Multiple Object Tracking Precision):} Określa geometryczną precyzję dopasowania przestrzennego ramek ograniczających dla poprawnie skojarzonych par obiektów~\cite{bernardin2008}:
        \begin{equation}
          MOTP = \frac{\sum_{i,t}d_{i,t}}{\sum_{t}c_{t}}
        \end{equation}
        gdzie $d_{i,t}$ reprezentuje odległość przestrzenną (wyrażoną najczęściej jako stopień pokrycia IoU) dla dopasowanej pary obiekt-detekcja $i$ w klatce $t$, a $c_t$ określa liczbę wszystkich poprawnych dopasowań w klatce $t$.
\end{itemize}

W celu wyeliminowania dominacji błędów detekcji nad asocjacją wprowadzono metryki zorientowane na
tożsamość (\textit{ang. ID-metrics})~\cite{ristani2016}:

\begin{itemize}
  \item \textbf{IDF1 (Identification F1 Score):} Mierzy zdolność systemu do długoterminowego utrzymywania spójnych tożsamości wzdłuż całych trajektorii, stanowiąc średnią harmoniczną precyzji identyfikacji ($IDP$) oraz czułości identyfikacji ($IDR$)~\cite{ristani2016}:
        \begin{equation}
          IDF1 = \frac{2\cdot{IDTP}}{2\cdot{IDTP} + IDFP + IDFN}
        \end{equation}
        Wskaźnik ten przypisuje najwyższą wagę ciągłości trajektorii, dzięki czemu  odzwierciedla skuteczność działania modułów Re-ID.
\end{itemize}

Współczesnym, najbardziej kompleksowym standardem oceny jest metryka \textbf{HOTA}, która
hierarchicznie balansuje dokładność lokalizacji i asocjacji:

\begin{itemize}
  \item \textbf{HOTA (Hierarchical Object Tracking Accuracy):} Definiowana jako średnia geometryczna z dokładności detekcji ($DetA$) oraz dokładności asocjacji ($AssA$), obliczana dla różnych progów IoU ($\alpha$)~\cite{luiten2021hota}:
        \begin{equation}
          HOTA = \sqrt{DetA \cdot AssA}
        \end{equation}
        gdzie składowe te reprezentują odpowiednio:
        \begin{equation}
          DetA = \frac{TP}{TP + FN + FP}, \quad AssA = \frac{TPA}{TPA + FNA + FPA}
        \end{equation}
        Wartość $TPA$ (\textit{ang. True Positive Association}) reprezentuje zbiór poprawnych przypisań tożsamości wzdłuż całej długości danej trajektorii, co pozwala na sprawiedliwą i nieobciążoną ocenę obu warstw systemu MOT jednocześnie.
\end{itemize}

\subsubsection{Wymagania i metryki wydajności sprzętowej}

W systemach realizujących analizę w czasie rzeczywistym, obok skuteczności matematycznej, kluczowa
jest ocena obciążenia sprzętowego oraz przepustowości potoku przetwarzania:

\begin{itemize}
  \item \textbf{Czas działania (Inference Time / Latency):} Definiuje sumaryczne opóźnienie obliczeniowe $\tau_{total}$ potrzebne do pełnego obsłużenia pojedynczej klatki wideo:
        \begin{equation}
          \tau_{total} = \tau_{det} + \tau_{assoc} + \tau_{reid} \quad [ms]
        \end{equation}
        gdzie $\tau_{det}$ to czas detekcji, $\tau_{assoc}$ to czas asocjacji lokalnej (np. BoT-SORT), a $\tau_{reid}$ to czas ekstrakcji wektora cech przez model transformera.

  \item \textbf{FPS (Frames Per Second):} Określa rzeczywistą przepustowość systemu wyrażoną liczbą klatek przetworzonych w ciągu sekundy:
        \begin{equation}
          FPS = \frac{1000}{\tau_{total}}
        \end{equation}
        Dla zapewnienia stabilnego działania w czasie rzeczywistym, system musi stale utrzymywać wartość parametru FPS na poziomie nie niższym niż częstotliwość rejestrowania strumieni wejściowych z kamer CCTV (typowa wartość to $\ge 25$~FPS).
\end{itemize}

\subsection{Standardowe benchmarki (MOTChallenge)}
Ocena porównawcza algorytmów śledzenia opiera się na standaryzowanej platformie testowej
MOTChallenge:
\begin{itemize}
  \item \textbf{MOT16 / MOT17:} Zestaw składający się z 14 sekwencji wideo reprezentujących zróżnicowane warunki środowiskowe (sceny miejskie, zmienne oświetlenie, ruch pieszy)~\cite{milan2016}. Edycja MOT17 dostarcza wysoce precyzyjne adnotacje oraz dedykowany zestaw detekcji wygenerowanych przez trzy publiczne detektory (DPM, Faster R-CNN, SDP), co zapewnia uczciwość porównań metod~\cite{milan2016, zhang2022}.

  \item \textbf{MOT20:} Zbiór danych zaprojektowany specjalnie dla scen o skrajnie wysokim zagęszczeniu tłumu (np. rynki miejskie, hale koncertowe), charakteryzujący się występowaniem gęstych skupisk obiektów generujących permanentne okluzje~\cite{dendorfer2020}.
\end{itemize}

%%%%%%%%%%%%%%%%%%%%%%%%%%%%%%%%%%%%%%%%%%%%%%%%%%%%%%%%%%%
%%%%%%%%%%%%%%%%%%%%%%%%%%%%%%%%%%%%%%%%%%%%%%%%%%%%%%%%%%%
\section{Systemy wielokamerowe - MTMC}
\subsection{Wprowadzenie do problematyki wielokamerowego śledzenia obiektów}

Rozwój infrastruktury wizyjnej oraz skokowy wzrost mocy obliczeniowej systemów analitycznych
doprowadziły do ewolucji systemów nadzoru wideo z podejścia lokalnego do w pełni rozproszonego.
Klasyczne algorytmy śledzenia w pojedynczej kamerze (\textit{ang. Single-Camera Tracking} -- SCT),
choć osiągają obecnie zadowalającą dokładność dzięki wykorzystaniu zaawansowanych głębokich sieci
neuronowych, napotykają fundamentalne ograniczenia przestrzenne. Pojedyncza kamera posiada
ograniczone pole widzenia, w którym obiekty często ulegają długotrwałej okluzji (zasłonięciu), a
jej fizyczne usytuowanie nierzadko narzuca trudną perspektywę obserwacji. Z tego powodu
zainteresowanie badaczy i przemysłu skupia się obecnie na systemach śledzenia wielu obiektów w
sieci kamer, powszechnie określanych w literaturze akronimami MTMC (\textit{ang. Multi-Target
  Multi-Camera Tracking}) lub MCMT (\textit{ang. Multi-Camera Multi-Target})~\cite{ristani2016}.

Głównym celem systemów MTMC jest łączenie danych z wielu strumieni wideo, pozwalając na przypisanie
i utrzymanie unikalnego, globalnego identyfikatora \textit{(ang. Global ID)} dla każdego śledzonego
obiektu w całej sieci monitoringu~\cite{ristani2016, nvidia_mtmc}. Systemy te stanowią fundament
nowoczesnych, zautomatyzowanych przestrzeni handlowych, inteligentnych magazynów logistycznych oraz
platform zarządzania ruchem miejskim. Pozwalają one na odtworzenie ciągłych, spójnych trajektorii
obiektów w czasie i przestrzeni. Jest to niezbędne do analizy przepływu klientów, wczesnego
wykrywania anomalii oraz automatyzacji i optymalizacji procesów przemysłowych~\cite{nvidia_mtmc}.

Przestrzenie, w których wdraża się systemy MTMC, różnią się zakresem pokrycia wizualnego. Wyróżnia
się dwa podstawowe sposoby rozmieszczenia kamer, które wpływają na dobór metod
analizy~\cite{woo2024}. Pierwszy obejmuje sieci kamer o pokrywających się polach widzenia
(\textit{ang. overlapping FOV}), dzięki którym ten sam obiekt może być obserwowany równocześnie z
różnych perspektyw. Ułatwia to łączenie informacji geometrycznych. Drugi scenariusz, częstszy w
praktyce i trudniejszy obliczeniowo, obejmuje układy z rozłącznymi polami widzenia (\textit{ang.
  non-overlapping FOV}). W takim przypadku obiekt może zniknąć z obrazu jednej kamery, przejść przez
martwą strefę (\textit{ang. blind spot}), a następnie pojawić się w innej, oddalonej kamerze po
czasie trudnym do oszacowania~\cite{ristani2016}. Oznacza to, że problem nie sprowadza się już
tylko do śledzenia ruchu, lecz staje się zadaniem reidentyfikacji wymagającym modelowania
zależności czasowo-przestrzennych~\cite{ristani2016}.

\subsection{Standardowa architektura potoku przetwarzania danych MTMC}

Złożoność obliczeniowa i~strukturalna śledzenia wielokamerowego wymaga zastosowania architektury
zorientowanej na mikroserwisy oraz kaskadowe przetwarzanie potokowe (\textit{ang. pipeline}).
Współczesne rozwiązania przemysłowe i~akademickie (m.in. platforma NVIDIA DeepStream SDK czy
architektury z~konkursów AI City Challenge) opierają się na powszechnie przyjętym paradygmacie
śledzenia opartego na detekcji (\textit{ang. tracking-by-detection})~\cite{nvidia_mtmc, riu2024}.
Paradygmat ten dzieli złożony problem na dwa główne, sekwencyjne etapy: lokalne śledzenie wewnątrz
poszczególnych kamer oraz późniejszą, globalną asocjację między kamerami~\cite{woo2024}.

W~typowym, dwuetapowym podejściu do MTMC, proces rozpoczyna się od wygenerowania krótkich ścieżek
(\textit{ang. tracklets}) dla każdego obiektu rejestrowanego niezależnie przez każdą
z~kamer~\cite{woo2024, ristani2016}. W~pierwszym kroku moduły detekcji oparte na głębokich sieciach
splotowych (CNN) lokalizują sylwetki ludzi lub pojazdów na poszczególnych klatkach obrazu.
Następnie lokalny moduł śledzący (SCT), wykorzystując techniki estymacji ruchu (np. filtry
Kalmana), łączy te niezależne predykcje w~krótkoterminowe ścieżki wewnątrz pojedynczego
widoku~\cite{bewley2016, aharon2022}. Operacja ta ma na celu odfiltrowanie błędnych detekcji,
uzupełnienie luk w~trajektoriach powstałych na skutek krótkotrwałych okluzji oraz zagregowanie cech
wizualnych. Pozwala to na wygenerowanie znacznie stabilniejszego i~wiarygodniejszego opisu
tożsamości przed przekazaniem danych do algorytmów reidentyfikacji~\cite{ristani2016, riu2024}.

Wygenerowane w~pierwszym kroku krótkie ścieżki (\textit{ang. tracklets}) są następnie przekazywane
do głównego etapu systemu MTMC, nazywanego globalną asocjacją między kamerami (\textit{ang.
  Multi-Camera Association} -- MCA)~\cite{woo2024}. Zadaniem tego modułu jest powiązanie lokalnych
ścieżek, pochodzących z~różnych kamer i~różnych przedziałów czasowych, z~jedną spójną tożsamością
globalną. Same informacje wizualne, takie jak deskryptory Re-ID, często nie wystarczają, ponieważ
różne obiekty mogą wyglądać bardzo podobnie, na przykład pracownicy w~jednakowych uniformach albo
identyczne samochody na skrzyżowaniu. Dlatego ważną rolę odgrywają moduły, które sprawdzają
zgodność skojarzeń w~czasie i~przestrzeni~\cite{riu2024}. Moduły te odwzorowują pozycje obiektów na
wspólną płaszczyznę sceny, a~następnie wyznaczają takie skojarzenie w~grafie, które daje
najmniejszy koszt obliczony na podstawie odległości przestrzennej i~podobieństwa
wizualnego~\cite{ristani2016, riu2024}.

Ponieważ system musi jednocześnie przetwarzać wiele strumieni wideo w~czasie rzeczywistym lub
z~niewielkim opóźnieniem (\textit{ang. near real-time}), jego architektura jest zwykle silnie
równoległa. Zgodnie ze współczesnymi rozwiązaniami stosowanymi w~praktyce kolejne klatki trafiają
do modułów inferencyjnych korzystających z~akceleracji sprzętowej~\cite{nvidia_mtmc}. Następnie
uzyskane metadane są przesyłane przez rozproszone brokery komunikatów, takie jak Apache Kafka, do
systemów bazodanowych przechowujących globalny graf tożsamości. Architektura mikroserwisowa ułatwia
skalowanie całego systemu, na przykład z~wykorzystaniem orkiestracji kontenerów w~środowisku
Kubernetes~\cite{nvidia_mtmc}.

\subsection{Problem rzutowania geometrycznego}

Niezależne obrazy z~poszczególnych kamer (we współrzędnych pikselowych $u, v$) należy sprowadzić do
wspólnego układu odniesienia. Najczęściej stosowanym układem jest globalna metryczna płaszczyzna
terenu (\textit{ang. Bird's-Eye View} -- BEV). W~scenach o~niskim zagęszczeniu standardową praktyką
jest wyznaczenie dolnej krawędzi ramki ograniczającej jako fizycznego punktu styku obiektu
z~podłożem, a~następnie rzutowanie tego punktu na globalną mapę za pomocą macierzy homografii lub
modelu projekcyjnego kamery z~wykorzystaniem jej parametrów zewnętrznych.

Niestety, podejście to zawodzi w~przypadkach scen o~wysokim zagęszczeniu tłumu lub przy złożonych
okluzjach. Obiekty znajdujące się bliżej obiektywu zasłaniają dolne partie sylwetek obiektów
znajdujących się w~głębi sceny. Niemożliwe staje się więc precyzyjne wyznaczenie ich punktu styku
z~podłożem, a~opieranie rzutowania na zniekształconej predykcji prowadzi do znacznych przesunięć
wyestymowanej pozycji na globalnej płaszczyźnie BEV.

Rozwiązanie tego problemu w sieciach kamer o pokrywających się polach widzenia \textit{(ang.
  overlapping FOV)} zostało zaprezentowane w~2008 roku jako koncepcja probabilistycznej mapy
zajętości (\textit{ang. Probabilistic Occupancy Map} -- POM)~\cite{fleuret2008}.

\subsection{Matematyczne sformułowanie algorytmu POM}

Algorytm probabilistycznej mapy zajętości (\textit{ang. Probabilistic Occupancy Map} -- POM) służy
do szacowania prawdopodobieństwa obecności osób na zdyskretyzowanej płaszczyźnie terenu (np.
podłodze) na podstawie obrazów binarnych uzyskanych w~procesie odejmowania tła, w~sposób niejawny
radząc sobie z~problemem okluzji~\cite{fleuret2008, mandeljc2013tracking}. W~podejściu tym
wykorzystano probabilistyczny model generatywny obrazu oraz przybliżenie wariacyjne pola średniego
(\textit{ang. mean-field approximation}), aby rozwiązać zadanie estymacji prawdopodobieństwa
\textit{a posteriori} $P(X \mid B)$, gdzie $X$ opisuje wektor stanów mapy zajętości, a~$B$ zbiór
obrazów binarnych ze wszystkich kamer~\cite{fleuret2008}.

\paragraph{Zmienne i założenia modelu}

Płaszczyzna terenu podlega dyskretyzacji na siatkę $G$ komórek, z~których każda reprezentuje
potencjalną lokalizację osoby~\cite{fleuret2008}. Dla każdej komórki $k \in \{1,\dots,G\}$
definiuje się binarną zmienną losową $X_k \in \{0,1\}$, gdzie $X_k = 1$ oznacza obecność
człowieka~\cite{fleuret2008}. Całą mapę zajętości opisuje wektor stanów $X = (X_1,\dots,X_G)$.
Głównym celem algorytmu jest estymacja prawdopodobieństwa \textit{a posteriori} $P(X \mid B)$
względem obserwowanych obrazów binarnych $B = (B^1,\dots,B^C)$ pozyskanych z~$C$
kamer~\cite{fleuret2008, mandeljc2013tracking}.

W~celu uproszczenia obliczeń model zakłada niezależność \textit{a priori} pomiędzy komórkami, tzn.
$P(X) = \prod_{k=1}^{G} P(X_k)$, oraz przyjmuje stałe, niezwykle rzadkie prawdopodobieństwo
zajętości pojedynczej komórki, np. $\epsilon_k = P(X_k = 1) \ll 1$~\cite{fleuret2008}.

Dla każdej kamery $c \in \{1,\dots,C\}$ wyznaczany jest obraz binarny $B^c$, uzyskany w~procesie
wyodrębniania tła (maska pierwszego planu). Algorytm POM buduje dla każdej z~nich syntetyczny obraz
$A^c$ na podstawie aktualnej mapy zajętości $X$. W~pozycji każdej komórki $k$, dla której $X_k=1$,
rzutowany jest prostokąt $A^c_k$ przybliżający sylwetkę stojącej osoby, widzianej z~perspektywy
kamery $c$~\cite{fleuret2008}. Formalnie, wygenerowany obraz syntetyczny jest wynikiem pikselowej
sumy logicznej (\textit{ang. union}) tych sylwetek:
\[
  A^c = \bigoplus_{k=1}^{G} X_k A^c_k
\]

Przyjmuje się, że poszczególne widoki z~kamer są warunkowo niezależne, a~gęstość prawdopodobieństwa
obrazu obserwowanego zależy wyłącznie od odpowiadającego mu obrazu
syntetycznego~\cite{fleuret2008}:
\[
  P(B \mid X) = \prod_{c=1}^{C} P(B^c \mid A^c) = \prod_{c=1}^{C} \frac{1}{Z} \exp\!\big(-\Psi(B^c, A^c)\big)
\]
gdzie $Z$ to stała normalizacyjna, a~$\Psi(\cdot,\cdot)$ jest miarą zdefiniowanej ad hoc
„pseudoodległości” pomiędzy obrazami binarnymi.

Miarę $\Psi(B,A)$ zaprojektowano tak, aby nakładała karę za niezgodność pikseli między obrazem
obserwowanym $B$ a~syntetycznym $A$:
\[
  \Psi(B,A) = \frac{1}{\sigma} \frac{\big|B \otimes (1 - A) + (1 - B) \otimes A\big|}{|A|}
\]
gdzie $\otimes$ oznacza iloczyn pikselowy (część wspólną), $|A|$ sumę wartości pikseli obrazu $A$
(pole powierzchni), a~$\sigma$ to parametr skali modelujący jakość algorytmu wyodrębniania
tła~\cite{fleuret2008}. Intuicyjnie, licznik wzoru zlicza liczbę błędnych pikseli (fałszywe alarmy
i~pominięcia względem obrazu syntetycznego), natomiast mianownik normalizuje tę wielkość przez pole
powierzchni prostokąta, co zapobiega faworyzowaniu obiektów o~większych rozmiarach w~procesie
optymalizacji~\cite{fleuret2008}.

Ponieważ bezpośrednie wyznaczenie dokładnego rozkładu \textit{a posteriori} jest niemożliwe
obliczeniowo ze względu na wykładniczo rosnącą przestrzeń kombinacji mapy zajętości, autorzy
algorytmu POM zastosowali przybliżenie wariacyjne pola średniego (\textit{ang. mean-field
  approximation})~\cite{fleuret2008}. Proces ten polega na poszukiwaniu uproszczonego rozkładu
prawdopodobieństwa (zakładającego niezależność poszczególnych komórek siatki), który minimalizuje
dywergencję Kullbacka-Leiblera względem rzeczywistego rozkładu~\cite{fleuret2008}. Prowadzi to do
układu równań punktu stałego, rozwiązywanego w~sposób iteracyjny. W~celu zapewnienia wysokiej
wydajności, niezbędnej do przetwarzania sygnału wideo z~wielu kamer, w~praktycznej implementacji
wykorzystano strukturę obrazów całkowych. Pozwala to na błyskawiczne aktualizowanie
prawdopodobieństw obecności obiektów w~kolejnych iteracjach, aż do osiągnięcia zbieżności
modelu~\cite{fleuret2008}.

\paragraph{Interpretacja w kontekście wizji komputerowej}

Kluczową zaletą algorytmu POM jest jego zdolność do naturalnego i~skutecznego radzenia sobie ze
zjawiskiem okluzji. W~modelu generatywnym obecność obiektu w~jednej komórce „tłumaczy” brak
widocznych pikseli pierwszego planu dla potencjalnego obiektu znajdującego się tuż za nim, wzdłuż
tej samej osi projekcji kamery~\cite{fleuret2008}. Pozwala to modelowi na poprawne rozróżnianie
złożonych konfiguracji osób, zarówno w~sytuacjach silnych skrótów perspektywicznych, jak i~w
przypadku częściowej utraty informacji z~niektórych kamer~\cite{fleuret2008}. Ponadto,
wykorzystanie algorytmu POM jako początkowego etapu detekcji na mapie zajętości zapewnia wysoką
precyzję lokalizacji obiektów we współrzędnych świata~\cite{mandeljc2013tracking}. Stanowi to
podstawę stabilnego śledzenia globalnych trajektorii w~czasie~\cite{mandeljc2013tracking}.

\section{Reidentyfikacja obiektów - Re-ID}
\label{sec:reidentyfikacja_reid}

\subsection{Definicja i sformułowanie problemu reidentyfikacji}
\label{subsec:definicja_reid}

Reidentyfikacja obiektów (\textit{ang. person/vehicle re-identification} -- Re-ID) umożliwia
kojarzenie tożsamości celów ruchomych przemieszczających się pomiędzy fizycznie rozłącznymi
obszarami, które są obserwowane przez kamery o~niepokrywających się polach widzenia (\textit{ang.
  non-overlapping fields of view})~\cite{ristani2016, zheng2016person}.

Proces ten definiuje się jako zadanie wyszukiwania informacji wizualnej (\textit{ang. image
  retrieval}) w relacji $1:N$~\cite{zheng2016person, zheng2017wild}. Dane wejściowe dla algorytmu
Re-ID składają się z~dwóch podstawowych zbiorów:
\begin{itemize}
  \item \textbf{Zapytania} (\textit{ang. query} lub \textit{probe}) -- reprezentowanego przez wycięty kadr obrazu zawierający sylwetkę poszukiwanego obiektu
        (człowieka lub pojazdu). Kadr ten jest najczęściej generowany przez moduł detekcji lub śledzenia wewnątrz pojedynczej kamery~\cite{ristani2016, zheng2016person}.
  \item \textbf{Galerii} (\textit{ang. gallery}) -- wielkoskalowej bazy danych ($N$ elementów) zawierającej kadry obiektów referencyjnych
        zarejestrowanych przez inne kamery w~sieci, wśród których algorytm poszukuje dopasowań~\cite{zheng2016person}. W warunkach rzeczywistych galeria ta jest
        często generowana automatycznie przez detektory wizyjne, co nieuniknienie wprowadza do niej szum w postaci fałszywych detekcji~\cite{zheng2017wild}.
\end{itemize}

Skuteczne rozwiązanie tego problemu opiera się na dwóch podstawowych elementach~\cite{liao2015}.
Pierwszym jest wyznaczenie możliwie najlepiej rozróżniającej reprezentacji cech dla zapytania i
obrazów z galerii. Drugim jest nauka miary podobieństwa, czyli określenie, jak porównywać obiekty w
przestrzeni cech, aby uporządkować obrazy z galerii od najbardziej do najmniej podobnych względem
zapytania~\cite{liao2015, zheng2016person}.

Głównym wyzwaniem jest to, że model musi opierać się wyłącznie na cechach wyglądu zewnętrznego,
takich jak kolor i~struktura odzieży czy ogólna sylwetka~\cite{rao2023}. Wynika to z~faktu, że
niska rozdzielczość oraz skośne ustawienie kamer utrudniają wydobycie precyzyjnych cech
biometrycznych, takich jak geometria twarzy czy tekstura tęczówki~\cite{ristani2016,
  zheng2017wild}. Dodatkowym utrudnieniem są duże różnice między obrazami tej samej osoby,
spowodowane zmianami oświetlenia, perspektywy widzenia kamery, pozy obiektu oraz częstymi
okluzjami~\cite{liao2015, zheng2017wild}.

\subsection{Klasyczne podejścia oparte na cechach projektowanych ręcznie}
\label{subsec:hand_crafted_reid}

Początkowo badania nad reidentyfikacją osób koncentrowały się na ręcznym projektowaniu deskryptorów
wizualnych (\textit{ang. hand-crafted features}) odpornych na transformacje geometryczne
i~fotometryczne sceny oraz na optymalizacji klasycznych miar odległości (\textit{ang. metric
  learning})~\cite{zheng2016person}. Do najczęściej cytowanych rozwiązań tej rodziny należą:
\begin{itemize}
  \item \textbf{Histogramy koloru}: Reprezentujące globalny lub lokalny rozkład barw w~różnych przestrzeniach (RGB, HSV, YUV). Metody te charakteryzowały się prostotą obliczeniową, lecz wykazywały skrajną wrażliwość na zmiany oświetlenia i~różnice w temperaturze barwowej i~balansie bieli kamer.
  \item \textbf{SDALF} (\textit{ang. Symmetry-Driven Accumulation of Local Features}): Algorytm opiera się na lokalizacji osi symetrii i~asymetrii ludzkiego ciała, co pozwala na podział sylwetki na stabilne regiony (głowa, tułów, nogi). Cechy lokalne (takie jak tekstury oparte na filtrach Gabora i~histogramy kolorów) są akumulowane wzdłuż tych osi, co pozwala na uzyskanie reprezentacji odpornej na rotacje obiektu i~zmiany pozy~\cite{farenzena2010person}.
  \item \textbf{LOMO} (\textit{ang. Local Maximal Occurrence}): Zaawansowana metoda opracowana przez Liao i~in.~\cite{liao2015}. Aby uodpornić system na wahania jasności, obrazy przetwarzane są transformatą Retinexa, która symuluje percepcję ludzkiego oka, uwydatniając detale w cieniach i stabilizując barwy~\cite{liao2015}. Do ekstrakcji tekstur wykorzystano niezależny od skali intensywności i wysoce odporny na szum operator SILTP~\cite{liao2015}. Główną innowacją LOMO jest jednak proces analizy poziomych pasów obrazu za pomocą przesuwających się okien. Algorytm rejestruje maksymalną wartość wystąpienia danej cechy (np. wzoru lub koloru) wzdłuż całej linii poziomej~\cite{liao2015}. Taka agregacja sprawia, że ostateczny wektor cech jest niezwykle odporny na rotację obiektu (zmiany punktu widzenia kamery – \textit{ang. viewpoint invariance}), eliminując potrzebę dokładnego lokalizowania poszczególnych partii ciała~\cite{liao2015}.
\end{itemize}

\subsection{Uczenie reprezentacji metrycznych i sieci syjamskie}
\label{subsec:metric_learning_reid}

Zastosowanie głębokich sieci neuronowych do reidentyfikacji pozwoliło na zjednoczenie etapu
ekstrakcji cech. Najważniejszą rolę w tym procesie odgrywa głębokie uczenie metryk (\textit{ang.
  deep metric learning}) w~podejściu \textit{end-to-end}, dodatkowe, osobne etapy klasycznego uczenia
odległości stały się zbędne~\cite{hermans2017defense}.

Wczesne architektury głębokie bazowały na sieciach syjamskich (\textit{ang. siamese networks})
składających się z~dwóch lub więcej bliźniaczych gałęzi o~współdzielonych wagach, do których
wprowadzano pary obrazów~\cite{zheng2017wild}. Sieci te optymalizowano przy użyciu funkcji straty
kontrastowej (\textit{ang. contrastive loss}), która dąży do minimalizacji odległości między
cechami par obrazów o~tej samej tożsamości (pozytywnych) oraz maksymalizacji odległości par
o~różnych tożsamościach (negatywnych), tak aby oddalić je od siebie powyżej zdefiniowanego
marginesu $\alpha$~\cite{zheng2017wild}.

Powszechnie stosowaną alternatywą stała się funkcja straty trójkowej (\textit{ang. triplet loss}),
która operuje na zestawach złożonych z~próbki kotwiczącej ($x_a$), próbki pozytywnej ($x_p$) oraz
próbki negatywnej ($x_n$)~\cite{hermans2017defense}. Klasyczny triplet loss dąży do spełnienia
warunku:
\begin{equation}
  D(x_a, x_p) + m < D(x_a, x_n)
  \label{eq:triplet_loss}
\end{equation}
gdzie $D(\cdot, \cdot)$ oznacza euklidesową odległość w~przestrzeni cech, a~$m$ reprezentuje narzucony margines bezpieczeństwa~\cite{hermans2017defense}.

Hermans i~in.~\cite{hermans2017defense} wykazali, że efektywność tej straty w~Re-ID ulega
drastycznej poprawie przy zastosowaniu strategii próbkowania najtrudniejszych przykładów wewnątrz
mini-batcha (\textit{ang. batch-hard triplet loss}):
\begin{equation}
  \mathcal{L}_{BH} = \sum_{i=1}^{P} \sum_{a=1}^{K}
  \left[ m + \max_{p=1 \dots K} D(x_i^a, x_i^p) - \min_{\substack{j=1 \dots P \\ j \neq i \\ n=1 \dots K}} D(x_i^a, x_j^n) \right]_+
  \label{eq:batch_hard_triplet}
\end{equation}
gdzie $P$ oznacza liczbę unikalnych tożsamości w~batchu, $K$ to liczba obrazów na każdą tożsamość, $x_i^a$ oznacza $a$-tą próbkę kotwiczącą należącą do $i$-tej tożsamości, a~$[\cdot]_+$ reprezentuje funkcję $\max(0, \cdot)$~\cite{hermans2017defense}. Algorytm ten koncentruje proces optymalizacji wyłącznie na najmniej podobnych parach pozytywnych oraz najbardziej zbliżonych parach negatywnych, eliminując dominację łatwych przykładów~\cite{hermans2017defense}.

\subsection{Architektury splotowe (CNN) dedykowane do zadań Re-ID}
\label{subsec:cnn_reid}

Większość splotowych modeli reidentyfikacji osób wykorzystuje architekturę ResNet-50 jako sieć
bazową (\textit{ang.backbone}) ze względu na korzystny kompromis między dokładnością a kosztami
obliczeniowymi~\cite{ristani2016, zheng2016person}. Ze względu na specyfikę zadania opracowano
dedykowane paradygmaty uczenia reprezentacji:
\begin{itemize}
  \item \textbf{IDE} (\textit{ang. ID-discriminative Embedding}): Model wprowadzony przez Zhenga i~in.~\cite{zheng2017wild}. Traktuje on proces treningowy jako problem klasyfikacji wieloklasowej, w którym każda tożsamość w~zbiorze uczącym stanowi odrębną klasę - model optymalizowany jest z użyciem funkcji straty opartej na entropii krzyżowej ~\cite{zheng2017wild}. Podczas wnioskowania odrzuca się warstwę klasyfikacyjną, a~wyjściowy wektor cech z~przedostatniej warstwy (GAP) służy bezpośrednio jako deskryptor tożsamości obiektu~\cite{zheng2017wild}.
  \item \textbf{PCB} (\textit{ang. Part-based Convolutional Baseline}): Model zaproponowany przez Suna i~in.~\cite{sun2018beyond} w~celu rozwiązania problemu utraty informacji przestrzennych w~globalnych deskryptorach. PCB rezygnuje z~globalnego uśredniania cech \textit{(GAP)} na rzecz równomiernego, sztywnego podziału wyjściowej mapy cech na $p$ poziomych pasów, najczęściej $p=6$~\cite{sun2018beyond}. W przeciwieństwie do wcześniejszych metod, podział ten ma charakter wyłącznie przestrzenny i~nie wymaga dopasowania do semantycznych części ciała. Dla każdego pasa obliczany jest niezależny lokalny wektor cech, optymalizowany osobną funkcją straty, co pozwala modelowi uchwycić drobne, lokalne szczegóły wyglądu (\textit{fine-grained features})~\cite{sun2018beyond}.
\end{itemize}

\subsection{Architektury oparte na Transformerach (SOTA)}
\label{subsec:transformers_reid}

Obecnie do najskuteczniejszych rozwiązań w dziedzinie Re-ID należą modele oparte na architekturze
Vision Transformer (ViT). Dzięki mechanizmowi samoatencji (\textit{ang. self-attention}) potrafią
one modelować zależności długiego zasięgu (\textit{ang. long-range dependencies}) bezpośrednio na
podstawie sekwencji fragmentów obrazu (\textit{ang. patches}), bez utraty szczegółów wynikającej z
operacji zmniejszania rozdzielczości (\textit{ang .downsampling}) charakterystycznych dla
klasycznych sieci splotowych~\cite{he2021transreid}. Przełomową architekturą tej grupy jest model
\textbf{TransReID} opracowany przez He i~in.~\cite{he2021transreid}. Wprowadza on do struktury
transformera dwa istotne moduły:
\begin{itemize}
  \item \textbf{JPM} (\textit{ang. Jigsaw Patch Module}): Moduł ten dzieli wektory cech fragmentów obrazu na kilka grup, dokonując na nich operacji przesunięcia (\textit{ang. shift}) i~przemieszania (\textit{ang. patch shuffle}) przed przekazaniem ich do kolejnych warstw sieci~\cite{he2021transreid}. Wymusza to na modelu uwzględnianie szerszego kontekstu globalnego i ogranicza jego nadmierne skupienie na jednym dominującym fragmencie sylwetki. W rezultacie zwiększa się odporność systemu na częściowe okluzje oraz zmiany pozy śledzonego obiektu~\cite{he2021transreid}.
  \item \textbf{SIE} (\textit{ang. Side Information Embeddings}): Pozwala na bezpośrednią integrację niewizualnych informacji kontekstowych, takich jak identyfikator kamery czy punkt widzenia, za pomocą wyuczalnych reprezentacji (\textit{ang. learnable embeddings})~\cite{he2021transreid}. Modyfikuje to wejściową sekwencję $Z_0$ zawierającą już informacje o~położenie zgodnie ze wzorem:
        \begin{equation}
          Z'_0 = Z_0 + \lambda S_{(C,V)}
          \label{eq:sie}
        \end{equation}
        gdzie $S_{(C,V)}$ oznacza połączoną reprezentację dla danej kamery i~perspektywy, natomiast~$\lambda$ jest parametrem wagowym. Mechanizm ten pozwala ograniczyć błędy w~przestrzeni cech wynikające ze specyfiki fizycznego położenia danej kamery, a tym samum zmniejszyć różnice w~podobieństwie obiektów rejestrowanych z~różnych ujęć~\cite{he2021transreid}.
\end{itemize}

\subsection{Problem rozbieżności dziedzin i adaptacja bez nadzoru (UDA)}
\label{subsec:domain_gap_reid}

Jednym z~najważniejszych ograniczeń praktycznych systemów Re-ID jest rozbieżność między
środowiskiem treningowym a~warunkami pracy systemu, czyli tzw. przesunięcie dziedzinowe
(\textit{ang. domain shift / domain gap})~\cite{ristani2016}. Model, który działa bardzo dobrze na
danych z~jednego zbioru kamer, po przeniesieniu do innej sieci monitoringu często wyraźnie traci
skuteczność~\cite{ristani2016, ge2020mutual}. Wynika to z~różnic w~obrazie rejestrowanym przez
różne kamery: inne jest oświetlenie, tło, perspektywa, jakość kompresji oraz sam wygląd osób
widzianych z~różnych punktów obserwacji~\cite{ristani2016, ge2020mutual}.

W~odróżnieniu od klasycznych problemów nadzorowanego uczenia, w~praktycznych zastosowaniach zwykle
nie ma możliwości ręcznego opisania dużej liczby danych z~nowego środowiska. Z~tego powodu
rozwijane są metody nienadzorowanej adaptacji dziedzinowej (\textit{ang. Unsupervised Domain
  Adaptation} -- UDA)~\cite{ge2020mutual}. Ich celem jest przeniesienie wiedzy z~domeny źródłowej do
docelowej bez użycia etykiet dla danych docelowych. W~literaturze Re-ID szczególnie ważne stały się
podejścia oparte na schemacie nauczyciel--uczeń, ponieważ pozwalają wykorzystywać nieopisane dane
z~nowego środowiska w~sposób kontrolowany i~stosunkowo stabilny.

Jednym z~najbardziej znanych rozwiązań tego typu jest algorytm \textbf{MMT} (\textit{ang. Mutual
  Mean-Teacher}) zaproponowany przez Ge i~in.~\cite{ge2020mutual}. Metoda ta wykorzystuje dwie
równolegle uczone sieci oraz ich wersje nauczycielskie, których parametry nie są aktualizowane
przez propagację błędu, lecz jako wykładnicza średnia krocząca (\textit{ang. Exponential Moving
  Average} -- EMA) wag modeli uczniów~\cite{ge2020mutual}. Dzięki temu model nauczyciela dostarcza
bardziej stabilnych pseudoetykiet dla danych docelowych.

Kluczowym elementem MMT jest wzajemne uczenie się modeli na podstawie miękkich pseudoetykiet
(\textit{ang. soft pseudo-labels}) generowanych dla danych z~domeny docelowej~\cite{ge2020mutual}.
Zamiast przypisywać próbkę do jednej, sztywnej klasy, model otrzymuje rozkład prawdopodobieństwa,
który lepiej oddaje niepewność predykcji. Takie podejście zmniejsza ryzyko błędów wynikających
ze~zbyt pewnych, ale niepoprawnych przypisań i~ogranicza narastanie szumu w~kolejnych iteracjach
uczenia~\cite{ge2020mutual}.

W~kontekście Re-ID jest to szczególnie istotne, ponieważ ta sama osoba może wyglądać bardzo różnie
w~zależności od kamery, kierunku patrzenia, odległości od obiektywu czy warunków oświetleniowych.
Dodatkowo w~obrazie często pojawia się tło, które nie niesie informacji o~tożsamości, a~jedynie
utrudnia dopasowanie. Z~tego powodu metody UDA dla Re-ID zwykle starają się nie tylko ograniczyć
różnice między domenami, ale też wydobyć cechy najbardziej związane z~samą osobą, a~nie
z~otoczeniem.

\subsection{Metryki oceny jakości reidentyfikacji}

Ponieważ zadanie reidentyfikacji pieszego w~swej strukturze jest tożsame z~problemami wyszukiwania
informacji wizualnych (\textit{ang. image retrieval}), do oceny skuteczności modeli Re-ID stosuje
się wyspecjalizowane wskaźniki odzwierciedlające jakość generowanego rankingu
podobieństwa~\cite{hermans2017defense, zheng2017wild}. Podstawowym standardem ewaluacji jest krzywa
skumulowanej charakterystyki dopasowania CMC (\textit{ang. Cumulative Matching Characteristics})
oraz średnia precyzja średnia mAP (\textit{ang. mean Average Precision})~\cite{ge2020mutual,
  zheng2017wild}.

Ewaluacja odbywa się poprzez porównanie wektora cech obrazu zapytania $q$ (tzw. \textit{probe} lub
\textit{query}) z~wektorami cech obrazów zgromadzonych w~bazie porównawczej (\textit{ang.
  gallery})~\cite{zheng2017wild}. Wynikiem działania systemu jest lista obrazów z~bazy uporządkowana
rosnąco według odległości metrycznej (np. cosinusowej) od zapytania~\cite{zheng2017wild}.

\begin{itemize}
  \item \textbf{Krzywa CMC i dokładność Rank-$k$:} Wykres CMC reprezentuje prawdopodobieństwo, że poprawne dopasowanie (obraz przedstawiający tę samą osobę, co zapytanie) znajdzie się wśród $k$ pierwszych (najbardziej podobnych) pozycji w~wygenerowanym rankingu~\cite{ge2020mutual, zheng2017wild}. Dla pojedynczego zapytania $q$, funkcja trafienia $f(q, k)$ przyjmuje wartość 1, jeśli poprawna tożsamość znajduje się w~zbiorze $k$ najlepiej dopasowanych obrazów z~bazy, oraz 0~w~przeciwnym razie. Skumulowaną dokładność Rank-$k$ dla całego zbioru zapytań $Q$ definiuje się jako średnią arytmetyczną~\cite{zheng2017wild}:
        \begin{equation}
          \text{Rank-}k = \frac{1}{|Q|} \sum_{q \in Q} f(q, k)
        \end{equation}
        Kluczową rolę odgrywa wskaźnik \textbf{Rank-1}, który określa zdolność systemu do poprawnego rozpoznania tożsamości już przy pierwszej, najbardziej podobnej próbie. Metryka ta charakteryzuje skuteczność wyszukiwania „najłatwiejszych” (najbardziej oczywistych) dopasowań w~sieci kamer (np. przy zbliżonej ekspozycji i~perspektywie)~\cite{sun2018beyond}.

  \item \textbf{Średnia precyzja średnia ($mAP$):} W~rzeczywistych scenach monitoringu miejskiego dany pieszy pojawia się wielokrotnie pod różnymi kamerami, przez co baza porównawcza zawiera wiele poprawnych obrazów docelowych~\cite{zheng2017wild}. Sama metryka Rank-1 jest wówczas niewystarczająca, ponieważ nie uwzględnia, czy system odnalazł pozostałe wystąpienia obiektu~\cite{sun2018beyond}. Problem ten rozwiązuje wskaźnik mAP, oceniający precyzję uporządkowania wszystkich prawidłowych trafień~\cite{zheng2017wild}.

        Dla pojedynczego zapytania $q$ wskaźnik średniej precyzji AP (\textit{ang. Average Precision})
        oblicza się według wzoru~\cite{zheng2017wild}:
        \begin{equation}
          AP(q) = \frac{\sum_{i=1}^{M} \left( P(i) \cdot \delta(i) \right)}{N_{gt}}
        \end{equation}
        gdzie $M$ oznacza całkowity rozmiar bazy porównawczej, $P(i)$ to precyzja wyznaczona dla top-$i$ zwróconych wyników, $\delta(i)$ jest funkcją indykatorową przyjmującą wartość 1, gdy $i$-ty obraz w~rankingu reprezentuje poszukiwaną tożsamość (oraz 0~w~przeciwnym razie), a~$N_{gt}$ określa sumaryczną liczbę obrazów danej osoby w~bazie (\textit{ground truth})~\cite{zheng2017wild}.

        Ostateczną wartość metryki $mAP$ dla całego eksperymentu oblicza się jako średnią z~wyników AP dla
        wszystkich zapytań ze zbioru $Q$~\cite{zheng2017wild}:
        \begin{equation}
          mAP = \frac{1}{|Q|} \sum_{q \in Q} AP(q)
        \end{equation}
\end{itemize}

Zastosowanie metryki mAP pozwala na rzetelną ocenę odporności modeli na trudne warunki środowiskowe
(silne zmiany oświetlenia i~pozy), gdyż wysoka wartość tego wskaźnika gwarantuje, że system jest
w~stanie odnaleźć wszystkie, również te mocno zdeformowane i~trudne wizualnie, wizerunki śledzonego
pieszego~\cite{sun2018beyond}.

\subsection{Specyfika reidentyfikacji pojazdów w odniesieniu do osób}
\label{subsec:vehicle_vs_person_reid}

Choć reidentyfikacja pojazdów (\textit{ang. vehicle Re-ID}) i~reidentyfikacja osób opierają się na
podobnych podstawach, oba zadania różnią się pod względem cech obiektu i~trudności praktycznych:
\begin{itemize}
  \item Sylwetka człowieka może zmieniać się w~czasie, na przykład w~zależności od pozy ciała, natomiast
        pojazd ma zwykle stały kształt.
  \item Samochody produkowane seryjnie często wyglądają bardzo podobnie, nawet jeśli są to różne
        egzemplarze tego samego modelu. Z~tego powodu sieć musi szukać drobnych, wyróżniających cech,
        takich jak tablice rejestracyjne, naklejki, uszkodzenia karoserii czy przedmioty widoczne we
        wnętrzu pojazdu.
  \item Wygląd pojazdu zależy od punktu widzenia kamery. Przód, bok i~tył auta mogą wyglądać zupełnie
        inaczej, dlatego w~praktyce często stosuje się podejście wielowidokowe (\textit{ang. multi-view
          Re-ID}).
\end{itemize}

Do oceny algorytmów Re-ID pojazdów stosuje się dedykowane zbiory danych, z~których najważniejsze to
\textbf{VeRi-776} oraz \textbf{VehicleID}. Pierwszy z~nich zawiera blisko 50~000 zdjęć 776 pojazdów
zarejestrowanych przez 20 kamer, a~drugi ponad 220~000 obrazów. Zbiory te stanowią obecnie
podstawowy punkt odniesienia przy testowaniu metod w~trudnych warunkach ruchu
drogowego~\cite{he2021transreid}.

\subsection{Prywatność i alternatywne metody identyfikacji}
\label{subsec:privacy_reid}

Tradycyjne podejścia do Re-ID, oparte na pełnych obrazach RGB sylwetek ludzi, budzą uzasadnione
obawy związane z~ochroną prywatności, zwłaszcza w~kontekście przepisów RODO. Z~tego powodu coraz
większe zainteresowanie budzą metody, które ograniczają ilość przetwarzanych danych już na etapie
projektowania systemu (\textit{ang. privacy by design}):
\begin{itemize}
  \item \textbf{Silhouette Re-ID}: Podejście polegające na identyfikacji osób wyłącznie na podstawie binarnych masek sylwetek wyodrębnionych z~obrazu, na przykład za pomocą algorytmów BGS. Metody te pomijają teksturę ubrań i~cechy twarzy, a~zamiast tego opierają się na kształcie sylwetki i~dynamice chodu (\textit{ang. gait-based Re-ID})~\cite{chao2019gaitset}. Powszechnie stosowany model \textbf{GaitSet} traktuje sekwencję chodu jako zbiór sylwetek, dzięki czemu dobrze radzi sobie z~brakującymi klatkami i~zmianami tempa poruszania się~\cite{chao2019gaitset}.
  \item \textbf{Skeleton-based Re-ID}: Podejście wykorzystujące estymację pozy, na przykład za pomocą OpenPose, do wyznaczenia współrzędnych kluczowych punktów ciała człowieka~\cite{rao2023}. Modele takie jak \textbf{TranSG} opisują ruch człowieka za pomocą grafu szkieletowego i~uczą się cech ruchu bez korzystania z~obrazu twarzy czy odzieży. Takie rozwiązanie zapewnia lepszą ochronę prywatności, a~jednocześnie pozwala zachować dobrą skuteczność rozpoznawania między kamerami~\cite{rao2023}.
\end{itemize}

\section{Rzutowanie na plan terenu (Bird’s-Eye View - BEV) i geometria sceny}

Widok z~lotu ptaka (\textit{ang. Bird's-Eye View} -- BEV) polega na rzutowaniu sceny na jedną,
wspólną płaszczyznę odniesienia, najczęściej na fizyczną płaszczyznę terenu lub jezdni.
Reprezentacja BEV pozwala na określenie położenia obiektów w~metrycznym układzie współrzędnych, co
całkowicie uniezależnia analizę od zniekształceń perspektywicznych specyficznych dla pojedynczej
kamery~\cite{fleuret2008, li2023}.

Współcześnie, w~systemach autonomicznej jazdy, BEV stanowi standardową przestrzeń reprezentacji
pośredniej (\textit{ang. intermediate representation}). Architektury percepcyjne przekształcają
surowe dane z~kamer lub sensorów LiDAR do widoku z~lotu ptaka, a~następnie dedykowane moduły
realizują w~tej zunifikowanej przestrzeni zadania takie jak detekcja obiektów 3D, semantyczna
segmentacja mapy oraz predykcja trajektorii ruchu~\cite{li2023, zhao2024}.

W~kontekście systemów monitoringu wizyjnego (CCTV), rzutowanie na płaszczyznę BEV pełni
analogiczną, fundamentalną rolę. Umożliwia ono geometryczną fuzję informacji z~wielu kamer w~ramach
jednego, globalnego układu współrzędnych, co stanowi podstawę do precyzyjnej analizy i~śledzenia
globalnych trajektorii obiektów w~czasie na planie terenu~\cite{fleuret2008}.

\subsection{Geometryczne rzutowanie na plan terenu i Inverse Perspective Mapping}

Klasycznym sposobem uzyskania widoku z~lotu ptaka jest rzutowanie obrazu perspektywicznego na
płaszczyznę terenu za pomocą macierzy homografii. W~podejściu tym zakłada się, że podłoże jest
idealnie płaskie, a~parametry kamery (macierz wewnętrzna oraz macierz rotacji i~wektor translacji)
są znane dzięki procesowi kalibracji~\cite{hartleyzisserman, szeliski2022}. Zgodnie z~modelem
kamery otworkowej, punkt w~przestrzeni trójwymiarowej rzutowany jest na płaszczyznę obrazu przy
użyciu macierzy projekcji $\mathbf{P} = \mathbf{K}[\mathbf{R} \mid \mathbf{t}]$. Zakładając, że
rozpatrywane punkty leżą bezpośrednio na płaszczyźnie gruntu (wysokość $Z=0$), przekształcenie to
redukuje się do macierzy homografii $\mathbf{H}$, która jednoznacznie odwzorowuje współrzędne
pikselowe obrazu $(u, v)$ na metryczne współrzędne w~układzie planu terenu $(x,
  y)$~\cite{hartleyzisserman, szeliski2022}.

Technika odwrotnego mapowania perspektywicznego (\textit{ang. Inverse Perspective Mapping} -- IPM)
polega na zastosowaniu tak wyznaczonej homografii do wszystkich pikseli obrazu
wejściowego~\cite{li2023}. W~rezultacie uzyskuje się syntetyczny obraz z~lotu ptaka, na którym
elementy infrastruktury drogowej (np. linie jezdni, pasy ruchu) oraz obiekty prezentowane są
z~pominięciem zniekształceń perspektywicznych~\cite{jeong2016}. IPM jest powszechnie stosowana
w~zadaniach związanych z~detekcją pasów i~oznakowania drogowego. Jej główną zaletą jest prostota
obliczeniowa oraz bezpośrednie oparcie na fizycznym modelu geometrycznym kamery. Fundamentalnym
ograniczeniem pozostaje jednak wysoka wrażliwość na błędy kalibracji (np. drgania pojazdu
zmieniające pochylenie kamery) oraz fakt, że założenie o~idealnie płaskim podłożu (\textit{ang.
  flat-ground assumption}) rzadko bywa w~pełni spełnione w~rzeczywistych warunkach
drogowych~\cite{jeong2016, li2023}.

\subsection{Siatka BEV i probabilistyczne mapy zajętości}

W~praktycznych systemach widok z~lotu ptaka jest zwykle reprezentowany jako siatka komórek
(\textit{ang. BEV grid}). Płaszczyzna terenu dzielona jest na prostokątne pola o~zadanym rozmiarze
(np. 10--50\,cm), z~których każda opisuje określony fragment
przestrzeni~\cite{mandeljc2013tracking}. Klasycznym podejściem, wywodzącym się z~robotyki mobilnej,
jest tworzenie siatek zajętości (\textit{ang. occupancy grid mapping}). Dla każdej komórki siatki
szacuje się prawdopodobieństwo, że jest ona zajęta przez przeszkodę lub obiekt. Informacje
z~różnorodnych czujników (kamery, LiDAR, radar) łączone są w~ramach probabilistycznego modelu,
a~mapa zajętości jest na bieżąco aktualizowana w~czasie~\cite{mandeljc2013tracking}.

Probabilistyczna mapa zajętości (\textit{ang. Probabilistic Occupancy Map} -- POM) zaproponowana
przez Fleureta i~in. jest w~istocie szczególnym przypadkiem siatki zajętości na płaszczyźnie
terenu~\cite{mandeljc2013tracking}. Autorzy dyskretyzują plan sceny na siatkę komórek i~estymują
prawdopodobieństwo obecności osoby w~każdym z~tych pól na podstawie obrazów binarnych uzyskanych
w~procesie wyodrębniania tła z~wielu kamer~\cite{fleuret2008}. W~niniejszej pracy siatka BEV jest
bezpośrednio utożsamiana z~siatką POM: dla każdej komórki $\mathcal{G}_k$ przechowywana jest
wartość brzegowego prawdopodobieństwa $P(X_k = 1)$, a~wynikowy widok z~lotu ptaka przedstawia
globalny rozkład zajętości obiektów na planie terenu w~danej chwili czasowej.

\subsection{Widok z lotu ptaka w systemach autonomicznej jazdy}

W badaniach nad autonomiczną jazdą widok z lotu ptaka (BEV) stał się jedną z kluczowych przestrzeni
reprezentacji pośredniej. Sieci neuronowe realizują transformację surowych danych z kamer i
sensorów LiDAR do przestrzeni BEV, a zadania takie jak detekcja trójwymiarowa, semantyczna
segmentacja mapy oraz predykcja ruchu pojazdów definiowane są bezpośrednio w tym układzie
współrzędnych~\cite{li2023, zhao2024}. Najnowsze przeglądy literaturowe podkreślają kilka
fundamentalnych zalet BEV: przejrzyste odwzorowanie sceny w układzie metrycznym, ułatwienie fuzji
danych z wielu różnorodnych sensorów, możliwość prostego definiowania stref na drodze (np. pasów,
skrzyżowań, obszarów wyłączonych z ruchu) oraz intuicyjne wykorzystanie tak zagregowanych
informacji przez moduły planowania ruchu i sterowania pojazdem~\cite{li2023, zhao2024}.

Te same cechy odgrywają istotną rolę w wielokamerowych systemach monitoringu wizyjnego (CCTV).
Widok z lotu ptaka pozwala na mapowanie położenia pieszych i pojazdów na punkty odpowiadające
rzeczywistym lokalizacjom w przestrzeni fizycznej, całkowicie uniezależniając analizę od
perspektywy narzucanej przez poszczególne kamery. Dzięki temu możliwa staje się analiza globalnego
ruchu na planie całej sceny, a nie tylko w obrębie odizolowanych obrazów
dwuwymiarowych~\cite{fleuret2008}.

\subsection{Metody BEV oparte na głębokim uczeniu}

Nowoczesne metody generowania widoku z~lotu ptaka odchodzą od bezpośredniego polegania na sztywnej
geometrii na rzecz modeli opartych na głębokich sieciach neuronowych. Uczą się one transformacji
z~widoku perspektywicznego do przestrzeni BEV w~sposób czysto oparty na danych, bez konieczności
ręcznego definiowania wszystkich kroków rzutowania~\cite{li2023}.

Przełomowym przykładem metody opartej na estymacji głębi (podejście 2D-3D) jest architektura
\textbf{Lift, Splat, Shoot (LSS)}~\cite{philion2020}. W~fazie \textit{Lift}, model estymuje
dyskretny rozkład głębi dla każdego piksela obrazu, a~następnie „podnosi” dwuwymiarowe cechy do
postaci trójwymiarowego ostrosłupa widzenia kamery (\textit{ang. viewing frustum})~\cite{zhao2024}.
Następnie, w~fazie \textit{Splat}, cechy te są rzutowane i~agregowane na płaską siatkę BEV. Z~kolei
faza \textit{Shoot} polega na wyznaczeniu z~tej reprezentacji mapy kosztów (\textit{ang. cost
  map}), służącej do semantycznej segmentacji sceny oraz planowania bezkolizyjnych trajektorii ruchu
pojazdu~\cite{philion2020, zhao2024}.

Alternatywnym, wysoce skutecznym podejściem jest model \textbf{BEVFormer}. Wykorzystuje on
architekturę Transformerów do bezpośredniego generowania reprezentacji BEV z~sekwencji obrazów
wielokamerowych~\cite{li2022bevformer}. Model ten wprowadza siatkę wyuczalnych wektorów zapytań
(\textit{ang. grid-shaped BEV queries}), z~których każdy odpowiada konkretnej komórce w~przestrzeni
BEV. Stosując zaawansowane mechanizmy atencji przestrzennej (\textit{ang. spatial cross-attention})
oraz czasowej (\textit{ang. temporal self-attention}), sieć elastycznie łączy i~filtruje informacje
z~różnych ujęć kamer oraz wcześniejszych klatek wideo~\cite{li2022bevformer}.

W~przeciwieństwie do klasycznej metody IPM, rozwiązania oparte na głębokim uczeniu wykazują
znacznie mniejszą wrażliwość na błędy kalibracji czujników i~lepiej radzą sobie w~scenach,
w~których założenie o~idealnie płaskim podłożu nie jest spełnione~\cite{li2023}. Ich główną wadą
pozostaje jednak konieczność posiadania potężnych zbiorów danych z~dokładnymi adnotacjami
w~przestrzeni 3D lub BEV, a~także charakterystyczny dla modeli głębokich brak pełnej przejrzystości
i~interpretowalności procesu decyzyjnego~\cite{li2023, zhao2024}.

\subsection{Zastosowanie widoku z góry w systemach wielokamerowych MTMC}

W~systemach śledzenia wielu obiektów w~sieci kamer (MTMC), widok z~lotu ptaka (BEV) stanowi
naturalną, wspólną przestrzeń reprezentacji. Rzutując detekcje z~poszczególnych kamer na plan
terenu, trajektorie obiektów można przedstawić jako sekwencje punktów w~jednej, globalnej
i~metrycznej siatce. Ułatwia to późniejsze kojarzenie ścieżek między kamerami oraz modelowanie
prawdopodobieństwa przejść~\cite{fleuret2008}. Probabilistyczna mapa zajętości (POM) dostarcza
w~każdej chwili precyzyjnych informacji o~tym, które komórki siatki są zajęte. Na tej podstawie
algorytmy mogą definiować strefy wejścia i~wyjścia, analizować główne kierunki ruchu, a~także
prognozować czas pojawienia się obiektu w~polu widzenia kolejnej kamery~\cite{fleuret2008}.

Jak zostanie zaprezentowane w~dalszej części pracy, docelowy system MTMC wykorzystuje widok z~lotu
ptaka jako absolutny fundament fuzji danych. Moduły śledzenia lokalnego (SCT) generują dla każdej
kamery niezależne trajektorie w~przestrzeni obrazu dwuwymiarowego, które następnie są
transformowane do układu BEV. Z~kolei moduł globalnej asocjacji kamer (MCA) łączy ze sobą cechy
wizualne (deskryptory Re-ID) oraz informacje czasowo-przestrzenne w~zunifikowanej przestrzeni
z~lotu ptaka, aby skutecznie utrzymać spójne identyfikatory obiektów w~całej monitorowanej sieci
monitoringu~\cite{woo2024, riu2024}.

\chapter{Przegląd istniejących rozwiązań}\label{chap:przegląd_rozwiązań}

W klasycznym ujęciu, architektura potoku przetwarzania danych MTMC dzieli się na dwa główne etapy:
lokalne śledzenie wewnątrz poszczególnych kamer (ang. Single-Camera Tracking – SCT), generujące
krótkie trajektorie (ang. tracklets), oraz późniejszą, globalną asocjację między kamerami (ang.
Multi-Camera Association – MCA), w której informacje wizualne i czasowo-przestrzenne są integrowane
w celu nadania unikalnych identyfikatorów. Wdrożenie tak zaawansowanego potoku na urządzeniach
klasy brzegowej lub stacjach roboczych o ograniczonej charakterystyce energetycznej wymusza
specyficzne, rygorystyczne podejście do doboru technologii oprogramowania. Niniejszy rozdział
stanowi przegląd istniejących rozwiązań algorytmicznych oraz architektonicznych, a także
uzasadnienie wyboru stosu technologicznego dla systemu MTMC budowanego w środowisku macOS. Docelową
platformą uruchomieniową jest komputer MacBook Pro wyposażony w układ Apple Silicon M5. Z uwagi na
unikalną specyfikę architektury sprzętowej Apple, standardowe i powszechnie stosowane w przemyśle
rozwiązania oparte na środowisku NVIDIA CUDA muszą zostać zastąpione alternatywnym stosem
optymalizacyjnym, gwarantującym najwyższą wydajność przy wykorzystaniu akceleratorów wbudowanych w
układ M5. Dokonano tu szczegółowej ewaluacji dostępnych języków programowania, modeli detekcji z
rodziny YOLO, głębokich architektur transformatorowych dostępnych za pośrednictwem platformy
Hugging Face oraz klasycznych algorytmów wizji komputerowej w ramach biblioteki OpenCV.

\section{Specyfika sprzętu oraz środowiska programistycznego}
\label{sec:ograniczenia_sprzętowe}

\subsection{Architektura sprzętowa: Apple Silicon M5 a akceleracja obliczeń tensorowych}

Głównym czynnikiem determinującym architekturę opisywanego systemu jest zastosowanie układu SoC
(\textit{ang. System-on-a-Chip}) z~rodziny Apple Silicon M5. Architektura ta wprowadza zmianę
w~sposobie zarządzania pamięcią operacyjną w~porównaniu do klasycznych środowisk opartych na
architekturze x86 i~dedykowanych kartach graficznych, co bezpośrednio wpływa na wydajność
przetwarzania modeli głębokiego uczenia~\cite{feng2025}.

Układy z~serii M wykorzystują zunifikowaną architekturę pamięci (\textit{ang. Unified Memory
  Architecture} -- UMA). Oznacza to, że procesor centralny (CPU), wielordzeniowy procesor graficzny
(GPU) oraz sprzętowy akcelerator macierzowy (\textit{ang. Neural Engine}) współdzielą tę samą,
wysoce przepustową pamięć RAM~\cite{feng2025, solovyova2026}. W~klasycznych systemach
wielokamerowych, opartych np. na systemie Ubuntu i~kartach z~serii NVIDIA RTX, konieczność ciągłego
przesyłania zdekodowanych klatek obrazu z~pamięci operacyjnej do pamięci karty graficznej (VRAM)
przez magistralę PCIe stanowi istotne ograniczenie wydajności potoku przetwarzania (\textit{ang.
  pipeline})~\cite{solovyova2026}.

Architektura UMA całkowicie eliminuje narzut czasu związany z~kopiowaniem tych danych. Układy
obliczeniowe uzyskują bezpośredni dostęp do surowych klatek wideo natychmiast po ich zdekodowaniu
przez procesor główny~\cite{solovyova2026}. Zmniejszenie tych opóźnień znacznie przyspiesza zarówno
procesy wstępnego przetwarzania obrazu, jak i~samo wnioskowanie (\textit{ang. inference}) na
modelach wizyjnych~\cite{solovyova2026}.

Układy Apple Silicon posiadają również sprzętowy akcelerator \textit{(ang. Neural Engine)},
zoptymalizowany pod kątem operacji mnożenia macierzy (\textit{ang. Matrix Multiplication} --
MatMul), które stanowią fundament obliczeniowy głębokich sieci neuronowych~\cite{solovyova2026}.
Aby w pełni wykorzystać ten potencjał, firma Apple opracowała interfejs programistyczny
\textit{(ang. Metal Performance Shaders -- MPS)}~\cite{feng2025}. Działa on jako niskopoziomowa
alternatywa dla przemysłowego standardu NVIDIA CUDA, bezpośrednio mapując operacje uczenia
maszynowego na architekturę układów z~serii M~\cite{feng2025, aliyev2025}.

Należy jednak zaznaczyć, że wsparcie popularnych frameworków (m.in. \textit{PyTorch}) dla
środowiska MPS wciąż znajduje się w~fazie rozwoju i~posiada pewne braki
optymalizacyjne~\cite{feng2025, aliyev2025}. Brak natywnej implementacji specyficznych operatorów
matematycznych w~API Metal skutkuje tzw. cichym przeniesieniem obliczeń na procesor główny
(\textit{ang. silent CPU fallback}). Może to powodować zatory w~potoku przetwarzania. Ograniczenie
to narzuca konieczność selekcji modeli i~wyboru wyłącznie takich architektur, których wszystkie
operacje są w~pełni wspierane sprzętowo przez GPU/Neural Engine. Jak zostanie to szerzej omówione
w~dalszej części pracy, uwarunkowanie to bezpośrednio zdeterminowało odrzucenie najnowszych
iteracji modeli YOLO w~niniejszej pracy.

\subsection{Środowisko programistyczne: Język Python oraz Miniforge}

Jako główną warstwę aplikacyjną systemu MTMC wybrano język Python, funkcjonujący w wirtualnym
środowisku Miniforge, będącym dystrybucją pakietów Conda dla architektury
ARM64~\cite{medium_miniforge}. Python stanowi obecnie standard w dziedzinie sztucznej inteligencji,
zapewniając dostęp do najszerszego ekosystemu bibliotek analitycznych. Mimo relatywnie niskiej
wydajności tego języka w~zadaniach obliczeniowych, w~kontekście nowoczesnych potoków głębokiego
uczenia pełni on funkcję wysokopoziomowego interfejsu (\textit{ang. wrapper}). Złożone obliczeniowo
operacje matematyczne są delegowane do zoptymalizowanych, bibliotek zaimplementowanych w~językach
C/C++ i wykonywane bezpośrednio na dedykowanych akceleratorach sprzętowych.

Oficjalne wdrożenie pełnego wsparcia backendu MPS do najpopularniejszego frameworka uczenia
maszynowego -- \textit{PyTorch}~\cite{apple_pytorch, feng2025}, uzasadnia wybór tego rozwiązania.
Dzięki temu ciężar operacji tensorowych może zostać przeniesiony bezpośrednio na układ graficzny
procesora M5, za pomocą zaledwie jednej instrukcji w~kodzie, definiując urządzenie sprzętowe jako
\texttt{device = torch.device("mps")}~\cite{apple_pytorch}. Bezpośrednia integracja omija
konieczność konwersji modeli do formatów pośrednich (np. CoreML) i~pozwala na płynną współpracę z
kluczowymi bibliotekami ekosystemu Python. Należą do nich m.in. NumPy - do wektoryzacji operacji w
zunifikowanej pamięci UMA, SciPy - do rozwiązywania problemów asocjacyjnych na grafach
dwudzielnych, oraz wbudowany moduł \texttt{multiprocessing}, wykorzystywany do izolowania procesów
dekodowania strumieni H.264 z~poszczególnych kamer.

\subsubsection{Alternatywne środowiska programistyczne}

Przez wzgląd na fakt, że system będzie działał na układach z~rodziny Apple Silicon, zrezygnowano
z~wykorzystania języka C++ jako głównego środowiska programistycznego. W~docelowych systemach
wbudowanych dla monitoringu wizyjnego C++ gwarantuje najwyższą kontrolę i~przewidywalność
zarządzania stertą (\textit{ang. heap}) oraz pamięcią podręczną, jednak na platformie macOS
wymuszałoby to implementację procesów wnioskowania poprzez niskopoziomowe interfejsy lub adaptację
lekkich silników brzegowych. Co więcej, implementacja najnowszych architektur opartych na
transformatorach -- które stanowią obecnie fundament nowoczesnych modułów reidentyfikacji
obiektów~\cite{he2021transreid} -- w~języku C++ i frameworku Metal znacznie wydłużyłaby czas
programowania, nie przynosząc przy tym współmiernych zysków dla całkowitej wydajności
systemu~\cite{feng2025}.

Z~analogicznych przyczyn, podyktowanych przede wszystkim brakiem nowoczesnych, oficjalnie
wspieranych przez społeczność dowiązań (\textit{ang. bindings}) do interfejsu MPS, odrzucono
również środowiska Java oraz MATLAB. Ponadto, jak uzasadniono we wcześniejszych podrozdziałach,
odrzucono architekturę NVIDIA CUDA oraz TensorRT ze względu na ich niekompatybilność z~układem SoC
M5~\cite{aliyev2025}.

\section{Śledzenie wewnątrz pojedynczej kamery i detekcja obiektów}

\subsection{Ekosystem YOLO oraz modele RT-DETR}

Do zadania lokalizacji obiektów w czasie niemal rzeczywistym wykorzystano jednoetapowe detektory
(\textit{ang. one-stage detectors}) zaimplementowane w bibliotece Ultralytics. Skupiono się przede
wszystkim na architekturze YOLOv8 oraz modelach opartych na transformatorach wizyjnych z rodziny
RT-DETR \textit{(ang. Real-Time Detection Transformer)}. Detektory jednoetapowe rozwiązują zadanie
detekcji poprzez bezpośrednią regresję współrzędnych ramek i klasyfikacji w jednym przebiegu sieci
(\textit{ang. forward pass}), co zapewnia im znaczną przewagę wydajnościową nad klasycznymi
podejściami opartymi na generowaniu regionów.

Z perspektywy wdrożenia systemu na układach Apple Silicon M5, kluczowym atutem biblioteki
Ultralytics jest pełna integracja akceleracji sprzętowej poprzez bezpośredni dostęp do interfejsu
Metal Performance Shaders (parametr \texttt{device='mps'}). Ważnym aspektem był jednak dobór
odpowiedniej generacji algorytmu. Mimo dostępności w literaturze nowszych wersji takich jak YOLOv10
lub YOLOv11, ostatecznie zdecydowano się na wykorzystanie YOLOv8, w wariantach Nano oraz Medium,
oraz modelu RT-DETR-L.

Decyzja ta wynika z dogłębnej analizy zachowania grafów obliczeniowych tych modeli na backendzie
MPS~\cite{ryspekov2026}. Architektura YOLOv8 opiera się na podejściu bezkotwicowym (\textit{ang.
  anchor-free}), rozdzielonej głowie detekcyjnej (\textit{ang. split head}) oraz regresji
dystrybucyjnej DFL (\textit{ang. Distribution Focal Loss}). Jej operacje splotowe wykazują pełną
zgodność z szablonami zoptymalizowanych warstw w Apple Metal. Z kolei nowsze architektury
wykorzystują eksperymentalne operacje atencji, których brak w bibliotece MPS skutkuje tzw. cichym
przeniesieniem obliczeń na powolny procesor główny (CPU)~\cite{ryspekov2026}. Dzięki optymalnej
kompatybilności, model YOLOv8-Nano uruchomiony na układach z serii Apple M wykazuje w testach
średnie opóźnienie rzędu 210~ms (ok. 4,75 FPS), przy jednoczesnym bardzo stabilnym odchyleniu
standardowym na poziomie 123~ms~\cite{ryspekov2026}.

Z kolei dla scen charakteryzujących się bardzo wysokim zagęszczeniem pieszych i długotrwałymi
okluzjami, klasyczne filtry bazujące na redukcji nadmiarowych ramek wykazują tendencję do błędnego
odrzucania pokrywających się sylwetek~\cite{ryspekov2026}. W takich sytuacjach system hybrydowo
integruje model RT-DETR-L. Formułuje on problem detekcji jako bezpośrednią predykcję zbioru
obiektów (\textit{ang. direct set prediction}), wykorzystując hybrydowy enkoder i algorytm
optymalnego skojarzenia w grafie dwudzielnym. Eliminacja heurystycznego algorytmu NMS zwalnia
znaczne zasoby procesora CPU~\cite{ryspekov2026}. Co więcej, dzięki optymalizacji operacji
macierzowych, model ten osiąga wysoką precyzję rzędu 54,71\% (mAP@0.5:0.95) przy stabilnym
opóźnieniu 331~ms~\cite{ryspekov2026}. Wybór ten stanowi więc kompromis pomiędzy precyzją niezbędną
do minimalizacji fragmentacji trajektorii, a opóźnieniem wymaganym w systemach bezpieczeństwa.

\subsection{Ograniczenia nowszych architektur YOLO w środowisku Apple MPS}

Choć w~środowiskach opartych na akceleratorach NVIDIA, modele YOLO v10 i v11 wykazują wyższą
precyzję przy mniejszej liczbie teoretycznych operacji zmiennoprzecinkowych, badania empiryczne
dowodzą, że wskaźniki te nie przekładają się na niskie opóźnienia na urządzeniach brzegowych
(\textit{ang. edge devices}), takich jak układy Apple Silicon~\cite{ryspekov2026}. Nowe generacje
YOLO, w~szczególności wersja v10, wprowadzają specyficzne moduły atencji (\textit{ang. Partial
  Self-Attention} -- PSA). Ponieważ framework Metal Performance Shaders wciąż posiada ograniczony
zakres wspieranych operatorów w~porównaniu ze środowiskiem CUDA, kompilator grafu obliczeniowego
nie jest w~stanie optymalnie zinterpretować bloków PSA~\cite{ryspekov2026}.

Z~powodu braku odpowiedniego wzorca sprzętowego, środowisko PyTorch wykonuje tzw. ciche
przeniesienie obliczeń (\textit{ang. silent CPU fallback}) z~wydajnych jednostek graficznych
z~powrotem na powolniejsze wątki procesora centralnego (CPU)~\cite{ryspekov2026, aliyev2025}.

Z~analogicznych względów wydajnościowych całkowicie zrezygnowano z~detektorów dwuetapowych, takich
jak architektura Faster R-CNN~\cite{ren2015faster}. Opierają się one na procesie generowania
propozycji regionów (RPN), którego iteracyjne przetwarzanie wyczerpałoby dopuszczalny budżet
opóźnień (\textit{ang. latency budget}) całego systemu. Z~kolei klasyczne podejścia z~dziedziny
wizji komputerowej (np. deskryptory HOG + SVM czy kaskady Viola-Jones) zostały odrzucone z~powodu
ich niskiej odporności na zniekształcenia perspektywiczne oraz braku zdolności adaptacyjnych do
luki domenowej (\textit{ang. domain gap}).

\subsection{Asocjacja danych i metody filtracji trajektorii: BoT-SORT i ByteTrack}

Najważniejszym kryterium doboru algorytmu asocjacji w pojedynczej kamerze był kompromis pomiędzy
odpornością na zjawisko okluzji a opóźnieniami w potoku przetwarzania. Choć klasyczny algorytm
SORT~\cite{bewley2016} cechuje się wydajnością rzędu 260~Hz, przez podatność na utratę tożsamości
obiektów w gęstych scenach nie mógł być zastosowany w systemie monitoringu~\cite{bewley2016}.
Wprowadzenie do procesu asocjacji cech wizualnych w architekturze Deep~SORT~\cite{wojke2017}
pozwoliło wprawdzie na redukcję liczby błędów zmiany tożsamości (\textit{ang. ID switches}) o 45\%,
jednak na współczesnych, zatłoczonych zbiorach danych model ten ustępuje precyzją nowszym
rozwiązaniom~\cite{wojke2017}.

Szczegółowej analizie porównawczej poddano architekturę ByteTrack~\cite{zhang2022} oraz jej
udoskonalone rozwinięcie -- BoT-SORT~\cite{aharon2022}. Średni czas wnioskowania ByteTrack wynosi
26,6~ms, jednak dla nietypowych warunków model ten traci precyzję lokalizacji
przestrzennej~\cite{singh2025}. Jak dowodzą najnowsze badania ewaluacyjne, dla obiektów szybko
poruszających się lub często zasłanianych, błędy śledzenia ByteTracka - wyrażone w średnim
przesunięciu względem celu - potrafią być o 260\% wyższe niż w konkurencyjnych metodach, osiągając
średnie odchylenie ADE na poziomie 114,3 piksela~\cite{singh2025}.

Z tego powodu ostatecznie zdecydowano się na wdrożenie algorytmu BoT-SORT~\cite{aharon2022}.
Rozwiązuje on problem wysokiego błędu przestrzennego i wykazuje znacznie wyższą stabilność
trajektorii przy wciąż akceptowalnym opóźnieniu. W testach na zbiorze MOT17, BoT-SORT osiągnął
najlepsze wyniki skuteczności (MOTA na poziomie 80,6 oraz IDF1 na poziomie 79,5), wyprzedzając tym
samym swojego poprzednika~\cite{aharon2022}.

Docelowa integracja niskowariancyjnych predykcji detektora YOLOv8 z asocjatorem BoT-SORT, który z
powodzeniem łączy w jednym kroku koszty geometryczne (IoU) i wizualne (Re-ID), gwarantuje rzetelną
bazę do utrzymania ciągłości trajektorii. Stanowi to warunek konieczny dla niezawodnego działania
globalnego modułu kojarzenia obiektów pomiędzy wieloma kamerami.

\section{Reidentyfikacja obiektów (Re-ID)}

\subsection{Vision Transformers i TransReID}

Ze względu na relatywnie niską rozdzielczość w systemach CCTV, która uniemożliwia identyfikację
biometryczną (np. na podstawie cech twarzy), rozpoznawanie obiektów musi opierać się wyłącznie na
elementach ubioru, geometrii oraz cechach szczególnych. Do implementacji modułu wybrano ekosystem
biblioteki \texttt{transformers} firmy Hugging Face, która posiada spójny interfejs API do
eksperymentów badawczych i jest kompatybilna z akceleratorami Apple Silicon. Transformacja modeli
odbywa się w kodzie w czasie rzeczywistym poprzez bezpośrednie lokowanie wag w architekturze Metal
(metoda \texttt{.to("mps")}).

Jako architekturę docelową wybrano rozwiązanie inspirowane modelem
TransReID~\cite{he2021transreid}, oparte o czyste transformatory wizyjne - ViT. Dzięki mechanizmowi
samoatencji (\textit{ang. self-attention}), architektury te wyodrębniają zależności przestrzenne
dalekiego zasięgu bezpośrednio z sekwencji wycinków obrazu (\textit{ang. patches}), zachowując przy
tym pełną rozdzielczość niezbędną do odróżniania wizualnie podobnych
obiektów~\cite{he2021transreid}. Zrezygnowano z implementacji klasycznych architektur opartych na
konwolucyjnych sieciach neuronowych (CNN), takich jak ResNet-50 ze schematem IDE czy popularny
model PCB~\cite{sun2018beyond}. Zasadniczą wadą sieci CNN w zadaniach Re-ID jest bezpowrotna utrata
drobnych, ale istotnych szczegółów sylwetki - logotypów na odzieży, krawędzi plecaków - wynikająca
z postępującej redukcji rozdzielczości w operacjach łączenia (\textit{ang. downsampling/pooling}).

\subsection{Wybór metod adaptacji domenowej i optymalizacji przestrzeni cech}

Wyzwaniem w monitoringu wielokamerowym pozostaje tzw. luka domenowa (\textit{ang. domain gap}),
objawiająca się drastycznymi różnicami w przestrzeni barw oraz perspektywie pomiędzy różnymi
punktami obserwacyjnymi. Aby zneutralizować ten problem i uodpornić system na zmienne warunki
środowiskowe, zdecydowano się na zastosowanie dwóch komplementarnych technologii.

Pierwszym zidentyfikowanym problemem były zniekształcenia geometryczne i perspektywiczne sylwetek.
Do ich korekty wybrano mechanizm bocznych osadzeń informacyjnych (\textit{ang. Side Information
  Embeddings - SIE})~\cite{he2021transreid}. Przez wzgląd na jego zdolność do bezpośredniego
przekazywania identyfikatorów poszczególnych kamer do wektora reprezentacji modelu. W
przeciwieństwie do klasycznych rozwiązań, SIE pozwala na jawną kompensację różnic wizualnych
wynikających z unikalnego umiejscowienia kamery, ułatwiając sieci rozpoznanie tego samego obiektu
pod różnymi kątami~\cite{he2021transreid}.

Drugim istotnym problemem projektowym był brak oznaczonych danych dla docelowego środowiska
działania systemu, a ręczne etykietowanie nowych nagrań jest procesem wysoce nieefektywnym. W celu
jego rozwiązania, proces optymalizacji oparto na architekturze nienadzorowanej adaptacji domenowej
(UDA). Spośród dostępnych metod wytypowano schemat MMT (\textit{ang. Mutual
  Mean-Teaching})~\cite{ge2020mutual}. MMT wykazuje odporność na zaszumione dane -- algorytm ten
poprawia zdolności uogólniania modelu poprzez generowanie miękkich pseudoetykiet przy użyciu
uśrednionych wag nauczyciela, skutecznie transferując wiedzę z domeny źródłowej do docelowej bez
jakiegokolwiek nadzoru eksperta~\cite{ge2020mutual}.

Do optymalizacji przestrzeni wyekstrahowanych cech na układzie Apple M5 wybrano funkcję straty
trójkowej z trudną eksploracją (\textit{ang. Batch-Hard Triplet Loss})~\cite{hermans2017defense}. W
przeciwieństwie do standardowych miar klasyfikacyjnych lub naiwnego doboru trójek, które prowadzą
do szybkiej stagnacji uczenia, metoda ta wymusza na sieci minimalizowanie odległości euklidesowej
$D$ wyłącznie dla najtrudniejszych przypadków pozytywnych i negatywnych w danej paczce:
\begin{equation}
  \mathcal{L}_{BH} = \sum_{i=1}^{P} \sum_{a=1}^{K} \left[ m + \overbrace{\max_{p=1 \dots K} D\left(f_\theta(x_i^{(a)}), f_\theta(x_i^{(p)})\right)}^{\text{najtrudniejszy pozytyw}} - \overbrace{\min_{\substack{j=1 \dots P \\ n=1 \dots K \\ j \neq i}} D\left(f_\theta(x_i^{(a)}), f_\theta(x_j^{(n)})\right)}^{\text{najtrudniejszy negatyw}} \right]_+
  \label{eq:batch_hard_triplet}
\end{equation}
gdzie $P$ oznacza liczbę unikalnych osób w paczce, $K$ liczbę obrazów danej osoby, a $m$ to założony margines separacji. Zastosowanie tej funkcji gwarantuje wyuczenie się wysoce dyskryminatywnych deskryptorów wizualnych, co bezpośrednio przekłada się na wyższą skuteczność kojarzenia trajektorii w globalnym module asocjacji~\cite{hermans2017defense}.

Ostatecznie, wdrożenie tak zaawansowanych mechanizmów musiało zostać dostosowane do specyfiki
sprzętowej procesora Apple M5. Ponieważ natywny backend MPS (\textit{Metal Performance Shaders})
posiada przejściowe braki w implementacji niektórych niskopoziomowych operatorów macierzowych,
architekturę uzupełniono o jawną deklarację zmiennej środowiskowej
\texttt{os.environ["PYTORCH\_ENABLE\_MPS\_FALLBACK"] = "1"}. Decyzja ta była niezbędna do
zabezpieczenia procesu ekstrakcji cech przed krytycznymi przerwaniami -- mechanizm ten wymusza
automatyczne i bezpieczne przeniesienie wykonania nieobsługiwanych fragmentów grafu obliczeniowego
na procesor centralny (CPU), gwarantując nieprzerwaną pracę całego systemu.

\subsection{Technologie odrzucone w zadaniu reidentyfikacji}

Podczas projektowania modułu Re-ID świadomie odrzucono metody oparte na ręcznie projektowanych
cechach wizualnych (\textit{ang. hand-crafted features}), w~tym zaawansowane deskryptory takie jak
LOMO (oparte na transformatach Retinex i~agregacji maksymalnych wystąpień w~przesuwnych
oknach)~\cite{liao2015} oraz moduły symetryczne SDALF~\cite{bazzani2013}. Choć charakteryzują się
one znacznie niższym narzutem obliczeniowym niż architektura TransReID, ich skuteczność drastycznie
spada w~przypadku zaszumionych nagrań o~niskiej rozdzielczości oraz przy znacznych zmianach
perspektywy (\textit{ang. viewpoint changes}) występujących pomiędzy różnymi kamerami. Zrezygnowano
również z~integracji klasycznych architektur syjamskich, trenowanych z~wykorzystaniem funkcji
straty kontrastowej w~starszych frameworkach (takich jak Caffe czy Theano). Stanowią one obecnie
przestarzałe rozwiązania, pozbawione możliwości płynnej adaptacji ich środowisk uruchomieniowych do
architektury zunifikowanej pamięci układów Apple M5.

Ostatecznie nie zdecydowano się także na wdrożenie metod realizujących założenia \textit{Privacy by
  Design}, czyli systemów bazujących na topologii szkieletowej (np. modele grafowe
TranSG~\cite{rao2023}) lub analizie dynamiki chodu (np. algorytm GaitSet~\cite{chao2019gaitset}).
Pomimo ich niewątpliwych zalet w~kontekście anonimizacji wizerunku, odrzucenie to podyktowane było
nieproporcjonalnie wysokim kosztem obliczeniowym niezbędnym do ekstrakcji punktów kluczowych
sylwetki. Wymagałoby to uruchomienia ciężkich estymatorów pozy (np. OpenPose) całkowicie
współbieżnie z~istniejącymi już detektorami obiektów, co natychmiast wyczerpałoby dostępny budżet
opóźnień całego systemu.

\section{Geometria sceny, rzutowanie BEV i modelowanie tła poprzez bibliotekę OpenCV}

\subsection{OpenCV z optymalizacjami dla architektury ARM}

Do realizacji podstawowych operacji na obrazach, dekodowania macierzy perspektywy oraz wstępnego
przetwarzania surowych strumieni wideo, wybrano powszechnie stosowaną bibliotekę OpenCV
(\textit{ang. Open Source Computer Vision}) w postaci interfejsu programistycznego dla języka
Python. Istotnym atutem tego rozwiązania w kontekście wdrożenia na układzie SoC Apple M5 jest
natywna optymalizacja jej bazowego kodu C++ pod procesory z rodziny ARM~\cite{opencv_changelog}.

Kluczowe algorytmy matematyczne OpenCV korzystają ze sprzętowych instrukcji wektorowych ARM NEON.
Odbywa się to poprzez tzw. uniwersalne instrukcje wewnętrzne (\textit{ang. universal intrinsics}),
deklarowane m.in. w nagłówku \texttt{arm\_neon.h}, które potrafią wektoryzować operacje w
środowisku zredukowanych rozkazów RISC na 128-bitowych rejestrach~\cite{opencv_changelog}. Pozwala
to na znaczną redukcję czasu wykonywania operacji zmiennoprzecinkowych (FP16) oraz
całkowitoliczbowych (SIMD) podczas analizy każdej z milionów klatek rejestrowanego wideo.
Dodatkowo, w środowisku macOS struktura OpenCV integruje się z niskopoziomowym pakietem Apple
Accelerate Framework. Jest to zbiór zoptymalizowanych procedur dostarczających funkcje numeryczne,
operacje algebry liniowej oraz biblioteki przetwarzania sygnałów. Przekłada się to na wysoką
wydajność podstawowych przekształceń, takich jak konwersje przestrzeni barw, rzutowania
perspektywiczne czy zmiany rozdzielczości, co skutecznie odciąża procesor główny, rezerwując pule
cykli obliczeniowych dla złożonych asocjacji logicznych w procesie śledzenia.

W celu skutecznego rzutowania detekcji pomiędzy wieloma kamerami, zwłaszcza w scenach o dużym
zagęszczeniu i silnych okluzjach, niezbędna jest transformacja geometrii wideo. W systemie
wykorzystano w tym celu Probabilistyczne Mapy Zajętości (POM, \textit{ang. Probabilistic Occupancy
  Map}), zaproponowane pierwotnie przez Fleureta i in.~\cite{fleuret2008}. Algorytm ten szacuje
prawdopodobieństwo zajętości danej komórki siatki na płaszczyźnie gruntu (reprezentowane zmienną
$X_k \in \{0,1\}$) przez pieszego. Estymacja ta opiera się bezpośrednio na maskach izolowanego tła
we wszystkich dostępnych kamerach~\cite{fleuret2008}. POM minimalizuje błędy rzutowania przy użyciu
probabilistycznego przybliżenia wariacyjnego, redukując odległość Kullbacka-Leiblera pomiędzy
estymowanym a rzeczywistym rozkładem prawdopodobieństwa obecności obiektów.

Zestawienie syntetycznego obrazu $A^c$, powstałego jako suma logiczna rzutowanych sylwetek ($A^c =
  \bigoplus_k X_k A^c_k$), z faktycznym obrazem masek tła $B$, pozwala na obliczenie kary (tzw.
pseudoodległości) za wszelkie niezgodności z obrazem rzeczywistym~\cite{fleuret2008}:
\begin{equation}
  \Psi(B,A) = \frac{1}{\sigma} \frac{\big|B \otimes (1 - A) + (1 - B) \otimes A\big|}{|A|}
\end{equation}
W powyższym równaniu operatory $\otimes$ oraz dodawania realizują operację różnicy symetrycznej (pomiędzy fałszywymi detekcjami a brakiem pokrycia), która następnie jest normalizowana przez powierzchnię $|A|$ i skalowana parametrem ufności $\sigma$.

Powyższy mechanizm, zwiększając odporność systemu na zjawisko wzajemnego przesłaniania się obiektów
wzdłuż osi optycznej kamery, skutecznie koryguje fizyczne współrzędne predykcji. Jako dane
wejściowe (zbiór $B$) do algorytmu POM podawane są maski pierwszoplanowe, uzyskane przy pomocy
technik wyodrębniania tła (\textit{ang. Background Subtraction} -- BGS) natywnie zaimplementowanych
w bibliotece OpenCV. W potoku przetwarzania wykorzystane zostały dwa podejścia. Pierwszym jest
wariant estymacji nieparametrycznej KNN, bazujący na estymatorze balonowym, który ogranicza
złożoność obliczeniową poprzez adaptację rozmiaru jądra do lokalnej gęstości
danych~\cite{zivkovic2006}. Drugim jest ulepszony algorytm mieszanin gaussowskich
MOG2~\cite{zivkovic2004}.

Zastosowanie algorytmu MOG2 skutecznie zwalnia system na sprzęcie brzegowym z wysokiego obciążenia
pamięciowego (\textit{ang. memory footprint}). Model ten przypisuje każdemu pikselowi obrazu
mieszaninę rozkładów normalnych z automatycznym kryterium optymalizacji ich liczby. Ostateczne wagi
w procesie mapowania MOG2 są aktualizowane z wykorzystaniem ujemnego współczynnika rozkładu a
priori Dirichleta:
\begin{equation}
  \hat{\pi}_m \leftarrow \hat{\pi}_m + \alpha (o_m^{(t)} - \hat{\pi}_m) - \alpha c_T
\end{equation}
Równanie to gwarantuje szybkie usunięcie przestarzałych klastrów tła (zjawisko tzw. eliminacji komponentów), które przestają otrzymywać wzmocnienie informacyjne, pozwalając modelowi na płynną adaptację do zmian w scenie~\cite{zivkovic2004}.

Całość operacji geometrycznych bazuje następnie na modelach kalibracyjnych i estymacji macierzy
homografii $3 \times 3$, które transformują punkty z obrazu perspektywicznego na płaszczyznę gruntu
za pomocą klasycznej metody odwrotnego mapowania perspektywicznego (IPM, \textit{ang. Inverse
  Perspective Mapping})~\cite{mallot1991}:
\begin{equation}
  \lambda \begin{bmatrix} x' \\ y' \\ 1 \end{bmatrix} = \mathbf{H} \begin{bmatrix} x \\ y \\ 1 \end{bmatrix}
\end{equation}
W środowisku OpenCV transformacja \texttt{cv::warpPerspective} jest bezpośrednio i natywnie wektoryzowana na sprzętowe jednostki SIMD (ARM NEON) w układzie M5, co pozwala na jej wykonanie w ułamku sekundy~\cite{opencv_changelog}.

Z rozważań projektowych definitywnie wykluczono współczesne algorytmy transformacji do przestrzeni
BEV oparte na głębokich sieciach neuronowych trenowanych w rygorze \textit{end-to-end} -- takie jak
modele Lift, Splat, Shoot (LSS)~\cite{philion2020} czy architektury bazujące na transformatorach
atencyjnych (np. BEVFormer~\cite{li2022bevformer}). Mimo że w domenie pojazdów autonomicznych
technologie te z powodzeniem wypierają modele analityczne (ignorując problem braku idealnej
płaskości terenu), ich wysokie wymagania obliczeniowe oraz zapotrzebowanie na ogromne zbiory danych
uczących czynią je niepraktycznymi w statycznym monitoringu wizyjnym (CCTV). Wymagałyby one
ciągłego ewaluowania wielkich, trójwymiarowych grafów wyłącznie po to, by wygenerować rzuty dla
sztywno zamontowanych kamer. Z tak sformułowanym problemem algorytmy analityczne OpenCV radzą sobie
natychmiastowo, korzystając ze stałych macierzy kalibracji. Odrzucono również starsze algorytmy
wyodrębniania tła, takie jak Bayesowski GMG, gdyż wymagają one długiej fazy inicjalizacyjnej
opóźniającej wejście systemu w tryb gotowości operacyjnej i wykazują niższą stabilność w
środowiskach wbudowanych~\cite{opencv_changelog}.

\chapter{Implementacja systemu śledzenia i reidentyfikacji}\label{chap:implementacja}

\section{Architektura potoku przetwarzania danych}

Wdrożone rozwiązanie analizy nagrań z systemu wielokamerowego opiera się na kaskadowym potoku
przetwarzania danych. Informacje wizualne są analizowane wieloetapowo: od lokalizacji
przestrzennej, poprzez śledzenie wewnątrz pojedynczej kamery, aż po globalną asocjację tożsamości.
Zastosowano architekturę współbieżną, system dystrybuuje zadania analizy wideo z pojedeynczej
kamery na niezależne podprocesy, następnie łącząc uzyskane dane w analizie widoku wielokamerowego.

\subsection{Środowisko uruchomieniowe i optymalizacja dla układu Apple M5}

Najważniejszym czynnikiem wpływającym na budowę prezentowanego systemu była specyfika platformy
sprzętowej – procesora z rodziny Apple Silicon M5. Aby wykorzystać potencjał obliczeniowy
wbudowanych akceleratorów macierzowych, główny framework głębokiego uczenia -- PyTorch został
zintegrowany z interfejsem MPS. Alokacja modeli neuronowych w pamięci koprocesora odbywa się
poprzez jawną deklarację urządzenia sprzętowego: \texttt{device = torch.device("mps")}. Ponieważ
interfejs API Metal nadal się rozwija i nie posiada natywnych implementacji dla wszystkich
zaawansowanych operatorów wykorzystywanych w architekturach Transformer, wprowadzono zabezpieczenie
na poziomie systemu operacyjnego. Poprzez deklarację zmiennej środowiskowej
\texttt{os.environ["PYTORCH\_ENABLE\_MPS\_FALLBACK"] = "1"} wymuszono automatyczne przekierowanie
nieobsługiwanych instrukcji na procesor główny, aby zapewnić bezawaryjną pracę potoku
przetwarzania.

\subsection{Wielowątkowość i asynchroniczna wymiana komunikatów}

Zarządzanie strumieniami wideo z poszczególnych kamer zrealizowano przy użyciu wbudowanej
biblioteki \texttt{multiprocessing}. Z uwagi na architekturę systemu macOS, do uruchamiania nowych
procesów wykorzystano metodę \texttt{spawn}, aby zapobiec konfliktom współdzielenia GPU podczas
duplikowania grafów obliczeniowych. Każda kamera obsługiwana jest przez niezależny proces roboczy
(\texttt{camera\_worker}). Proces ten samodzielnie dekoduje sygnał H.264, wykonuje detekcję oraz
ekstrahuje cechy wizualne. Cechą szczególną zaproponowanej architektury jest model komunikacji
międzyprocesowej. Do głównego procesu asocjacji przesyłane są wyłącznie ustrukturyzowane metadane w
postaci słowników w języku Python, zamiast klatek wideo. Pakiety te zawierają znormalizowane
wektory cech Re-ID, wyliczone współrzędne na płaszczyźnie oraz identyfikator kamery. Takie
podejście redukuje obciążenie magistrali pamięciowej, umożliwiając łatwą skalowalność systemu
poprzez dodawanie kolejnych sensorów wizyjnych.

\begin{figure}[H]
  \centering
  \includegraphics[width=.8\textwidth]{architektura.png}
  %\vspace*{-10mm}
  \caption{Wizualizacja architektury potoku przetwarzania danych. Źródło: opracowanie własne.}
  \label{fig:architektura}
\end{figure}

\section{Moduł detekcji i śledzenia w pojedynczej kamerze (SCT)}

Realizacja śledzenia wieloobiektowego opiera się na dominującym paradygmacie śledzenia przez
detekcję (\textit{ang. tracking-by-detection}). Zgodnie z tym podejściem problem rozdzielony jest
na lokalizowanie przestrzenne obiektów oraz późniejsze łączenie ich w ciągłe trajektorie.

\subsection{Inicjalizacja i wnioskowanie modelu YOLO}

Jako fundament modułu lokalizacji przestrzennej wybrano model z rodziny YOLOv8 w wariancie
\texttt{Medium} -- YOLOv8m. Przeprowadzono testy różnych modeli i zastosowany wariant oferował
najkorzystniejszy kompromis pomiędzy czasem wnioskowania a precyzją wykrywania częściowo
zasłoniętych sylwetek w zagęszczonych scenach miejskich.

Dla optymalizacji obciążenia logiki decyzyjnej procesu detekcji, ograniczono ilość danych
wejściowych. Próg ufności (\textit{confidence score}) sieci splotowej ustalono na poziomie $25\%$,
odrzucając predykcje o niskim stopniu pewności. Następnie nałożono filtr semantyczny, redukując
wyjściową przestrzeń wyników ze wszystkich $80$ klas zbioru COCO wyłącznie do klasy $0$,
odpowiadającej osobom. Wynikiem działania sieci jest wektor współrzędnych ramki otaczającej $b =
    [x_1, y_1, x_2, y_2]^T$, na podstawie którego wyznaczany jest centralny punkt dolnej krawędzi
reprezentujący fizyczny styk obiektu z podłożem.

\subsection{Śledzenie lokalne z wykorzystaniem algorytmu BoT-SORT}

Aby odfiltrować błędne detekcje oraz zniwelować efekty krótkotrwałych okluzji (np. przejście
pieszego za znakiem drogowym), wewnątrz każdego procesu kamery zaimplementowano lokalny algorytm
asocjacji BoT-SORT. Model ten bezpośrednio modeluje szerokość i wysokość ramki otaczającej w
strukturze filtru Kalmana.

Filtr Kalmana w każdej klatce estymuje przyszłe położenie obiektu na podstawie jego dotychczasowej
prędkości wektorowej. BoT-SORT wykorzystuje algorytm kompensacji ruchu do niwelowania ewentualnych
drgań obrazu, a następnie wyznacza macierz kosztów asocjacji. Działając jako bufor wygładzający,
śledzenie lokalne grupuje dyskretne detekcje pieszych w spójne fragmenty trajektorii (\textit{ang.
  tracklets}). Pozwala to na wysłanie do globalnego modułu wiarygodniejszej informacji o tożsamości,
odciążając tym samym wyższe warstwy systemu.

\begin{table}[htbp]
  \centering
  \caption{Parametry detektora i trackera lokalnego. Źródło: opracowanie własne.}
  \label{tab:params_local_tracker}
  \renewcommand{\arraystretch}{1.2}
  \begin{tabular}{|l|c|p{8cm}|}
    \hline
    \textbf{Parametr}                         & \textbf{Wartość} & \textbf{Opis}                                                                                    \\ \hline
    \textbf{Akceleracja sprzętowa}            & Apple MPS        & Wymuszenie wnioskowania na GPU procesora M5                                                      \\ \hline
    \textbf{Model detekcji}                   & YOLOv8m          & Wersja o optymalnym stosunku wydajności do precyzji                                              \\ \hline
    \textbf{Classes}                          & 0                & Ograniczenie detekcji do klasy pieszych                                                          \\ \hline
    \textbf{Próg ufności - \texttt{conf}}     & 0.25             & Minimalne prawdopodobieństwo uznania ramki otaczającej za reprezentującą rzeczywisty obiekt.     \\ \hline
    \textbf{Próg pokrycia - \texttt{iou}}     & 0.50             & Próg geometrycznego pokrycia ramek dla algorytmu asocjacji oraz redukcji duplikatów (NMS).       \\ \hline
    \textbf{Tracker lokalny}                  & BoT-SORT         & Zintegrowany algorytm śledzenia generujący krótkie trajektorie wewnątrz pojedynczej kamery.      \\ \hline
    \textbf{Wiek ścieżki - \texttt{max\_age}} & 30               & Maksymalna liczba klatek, przez które tracker podtrzymuje tożsamość ukrytego/zgubionego obiektu. \\ \hline

  \end{tabular}
\end{table}

\section{Ekstrakcja cech wizualnych i reidentyfikacja (Re-ID)}

Przez wzgląd na ograniczone pole widzenia pojedynczej kamery, niezbędny jest moduł reidentyfikacji,
odpowiedzialny za obliczanie trajektorii obiektów poruszających się między kamerami. Zadaniem tego
etapu jest wyznaczenie unikalnej sygnatury wizualnej, odpornej na zmiany perspektywy i warunków
oświetleniowych.

\subsection{Implementacja architektury Vision Transformer}

Czerpiąc z rozwiązań zaproponowanych w modelu TransReID, zaimplementowano rozwiązanie oparte o
trasnofrmatory wizyjne ViT. Mechanizm samoatencji (\textit{self-attention}) dzieli wejściowy
wycinek ramki pieszego na równe fragmenty (\textit{patches}), modelując zależności przestrzenne o
dalekim zasięgu pomiędzy nimi na przestrzeni wszystkich warstw, co znacznie zwiększa odporność na
częściowe zasłonięcia sylwetki.

\subsection{Generowanie i normalizacja wektorów cech (Embeddings)}

Wykorzystując biblioteki ekosystemu \texttt{Hugging Face}, wykadrowana sylwetka pieszego jest
formatowana do wymagań wejściowych modelu, rozmiar zmieniany jest do $224 \times 224$ pikseli, a
przestrzeń barw konwertowana z domyślnej dla OpenCV BGR na RGB. Proces wnioskowania
(\textit{forward pass}) realizowany jest w wysoce zoptymalizowanym kontekście
\texttt{torch.no\_grad()}, co zapobiega alokacji pamięci na obliczanie gradientów.

Rezultatem operacji jest tensor zagregowanej reprezentacji cech całej sylwetki (tzw.
\textit{embedding}) o wymiarach $1 \times 768$. Uzyskany wektor $\mathbf{v}$ poddawany jest na
wyjściu normalizacji L2 według wzoru: $$ \mathbf{v}_{norm} = \frac{\mathbf{v}}{||\mathbf{v}||_2} $$
Operacja ta nakłada wektory cech na jednostkową sferę wielowymiarową, co jest warunkiem koniecznym
do poprawnego zastosowania odległości cosinusowej w procesie globalnej oceny podobieństwa.

\section{Globalna asocjacja danych (Multi-Camera Association)}

Wielowątkowy potok danych zbiega się w obiekcie klasy \texttt{GlobalTracker}, zarządzającego
tożsamościami. Głównym wyzwaniem tego modułu jest matematyczne udowodnienie, że sylwetki
wyekstrahowane z różnych kamer w różnym czasie należą do tego samego pieszego.

\subsection{Konstrukcja macierzy kosztów i łączenie błędów}

Algorytm w każdej klatce buduje globalną macierz odległości, oceniając stopień dopasowania
wszystkich nowo nadesłanych detekcji względem obiektów istniejących w bazie. Problem
reidentyfikacji rozszerzono o weryfikację fizyczną, tworząc hybrydową macierz kosztów, na którą
składają się:
\begin{enumerate}
  \item \textbf{Koszt wizualny:} Wyliczany przy pomocy zoptymalizowanej funkcji z biblioteki SciPy, polegający na wyznaczeniu odległości cosinusowej pomiędzy znormalizowanymi wektorami Re-ID. Odległość bliska zeru oznacza skrajnie wysoką zbieżność tekstur i proporcji ubioru.
  \item \textbf{Koszt przestrzenny:} Metryka euklidesowa wyznaczana na podstawie pozycji obiektów przeniesionych na globalną płaszczyznę mapy rzutu pionowego (Bird's-Eye View). Weryfikacja dokładnego procesu transformacji współrzędnych z ujęcia kamery na mapę BEV zostanie szczegółowo omówiona w rozdziale \ref{chap:integracja_mapy}.
\end{enumerate}

W celu redukcji obciążenia obliczeniowego wprowadzono zaporowe progi odcięcia (np.
\texttt{max\_visual\_cost}). Jeśli badana para przekracza ustalony dystans euklidesowy lub
cosinusowy, przypisywany jest jej ryczałtowo skrajnie wysoki koszt (tzw. koszt zaporowy), co
natychmiast eliminuje ją z grafu poszukiwań jako parę fizycznie niemożliwą.

\subsection{Dopasowanie z użyciem algorytmu węgierskiego}

Wygenerowana hybrydowa macierz kosztów stanowi wejście dla dyskretnej optymalizacji algorytmu
węgierskiego (metody Kuhna-Munkresa). Wykorzystano w tym celu implementację
\texttt{linear\_sum\_assignment} pakietu \texttt{scipy.optimize}. Algorytm ten rozwiązuje problem
asocjacji w grafie dwudzielnym w czasie wielomianowym, wyznaczając takie relacje "jeden do
jednego", które minimalizują globalną sumę przypisań dla całego skrzyżowania. Stanowi to kluczowy
element radzenia sobie z problemem martwych stref między polami widzenia, gwarantując najmniejszy
błąd asocjacji, nawet gdy w obrębie ujęć kilku kamer pojawią się wizualnie podobni przechodnie.

\subsection{Zarządzanie cyklem życia globalnych identyfikatorów}

Ostatnim etapem procesu jest zatwierdzenie wyników asocjacji oraz aktualizacja grafu tożsamości.
Pomyślnie dopasowane detekcje natychmiast odświeżają wektory cech oraz globalne pozycje fizyczne w
obiekcie \texttt{GlobalTracker}.

Zarządzanie trajektoriami jest silnie deterministyczne. Nowe wykrycia, dla których algorytm
węgierski nie zdołał dopasować istniejącej historii, rejestrowane są jako nowe globalne
identyfikatory (\textit{Global ID}). Równolegle zastosowano mechanizm obsługi tzw. wygasających
ścieżek. W sytuacji, gdy przechodzień trwale opuści pole widzenia wszystkich kamer monitoringu,
jego identyfikator nie jest od razu porzucany. System inkrementuje licznik brakujących klatek
(\texttt{missed\_frames}), przechowując tożsamość w stanie uśpienia. Dopiero po przekroczeniu
ustalonego bufora czasowego, struktury przechowujące jego wektory cech są bezpowrotnie zwalniane z
pamięci operacyjnej, zapobiegając przepełnieniu zunifikowanej pamięci RAM układu.

\usetikzlibrary{positioning}
\begin{figure}[htbp]
  \centering
  \begin{tikzpicture}[
      global_node/.style={circle, draw=blue!80, fill=blue!10, thick, minimum size=1.1cm, align=center},
      det_node/.style={circle, draw=orange!80, fill=orange!10, thick, minimum size=1.1cm, align=center},
      matched_edge/.style={->, >=stealth, ultra thick, green!50!black},
      rejected_edge/.style={->, >=stealth, thin, dashed, gray},
      font=\small
    ]

    % Lewa kolumna: Global IDs
    \node[global_node] (G1) at (0, 3) {$G_1$};
    \node[global_node] (G2) at (0, 0) {$G_2$};

    % Prawa kolumna: Detekcje z kamer
    \node[det_node] (D1) at (6, 3.5) {$D_{c1,1}$};
    \node[det_node] (D2) at (6, 0.5) {$D_{c2,1}$};
    \node[det_node] (D3) at (6, -2.5) {$D_{c2,2}$};

    % Podpisy kolumn
    \node[above=0.3cm of G1, font=\bfseries] {Istniejące tożsamości (Global ID)};
    \node[above=0.3cm of D1, font=\bfseries] {Nowe detekcje (Z kamer)};

    % Krawędzie od G1
    \draw[matched_edge] (G1) -- node[above, sloped, font=\footnotesize] {$C = 0.15$ (Przyjęte)} (D1);
    \draw[rejected_edge] (G1) -- node[above, near end, sloped, font=\tiny] {$C = 0.62$} (D2);

    % Krawędzie od G2
    \draw[rejected_edge] (G2) -- node[below, near start, sloped, font=\tiny] {$C = 0.78$} (D1);
    \draw[matched_edge] (G2) -- node[below, sloped, font=\footnotesize] {$C = 0.08$ (Przyjęte)} (D2);

    % Nowy obiekt bez dopasowania
    \node[right=0.2cm of D3, font=\footnotesize\itshape, color=red!70!black] {Brak połączenia ($C > C_{max}$) $\rightarrow$ Tworzenie $G_3$};

  \end{tikzpicture}
  \caption{Ilustracja procesu globalnej asocjacji danych za pomocą grafu dwudzielnego. Algorytm węgierski szuka takich połączeń (pogrubione zielone linie), które minimalizują sumę kosztów $C$ opartych na odległości Re-ID oraz geometrii BEV. Detekcja $D_{c2,2}$ nie spełnia progów podobieństwa i inicjuje nowy identyfikator $G_3$.}
  \label{fig:bipartite_matching}
\end{figure}

\chapter{Integracja widoku wielokamerowego i rzutowanie na mapę}\label{chap:integracja_mapy}

Zastosowanie pojedynczych kamer w systemach wizyjnych dostarcza precyzyjnych informacji o
lokalizacji obiektów wyłącznie w dwuwymiarowej przestrzeni obrazu (wyrażonej w pikselach).
Rozwiązanie problemu śledzenia obiektów w rozproszonej sieci kamer (MTMC) wymaga sprowadzenia tych
niezależnych obserwacji do wspólnego układu odniesienia. W niniejszym rozdziale przedstawiono
implementację transformacji perspektywicznej do układu rzutu pionowego, powszechnie określanego
jako widok z lotu ptaka (\textit{ang. Bird's-Eye View} -- BEV). Opisano również metody weryfikacji
przestrzennej oparte na probabilistycznej mapie zajętości (POM), które korygują błędy rzutowania
wynikające z okluzji w gęstym ruchu miejskim

\section{Wyznaczanie i aplikacja macierzy homografii}
Fundamentem integracji geometrycznej w zaimplementowanym systemie jest technika odwrotnego
mapowania perspektywicznego (\textit{ang. Inverse Perspective Mapping} -- IPM). Przyjmując
założenie, że obserwowany obszar stanowi płaską płaszczyznę terenu ($z=0$), pełny model projekcji
kamery redukuje się do płaskiego przekształcenia rzutowego, opisywanego przez macierz homografii
$\mathbf{H}$ o wymiarach $3 \times 3$.

\subsection{Model geometryczny i odwrotne mapowanie perspektywiczne (IPM)}

Zależność pomiędzy punktem na płaszczyźnie obrazu z kamery a jego rzeczywistym położeniem na
globalnej mapie BEV opisana jest równaniem: $$ \lambda \begin{bmatrix} x' \\ y' \\ 1 \end{bmatrix} = \mathbf{H} \begin{bmatrix} x \\ y \\ 1 \end{bmatrix} $$
gdzie $(x, y)$ to współrzędne pikselowe w obrazie źródłowym, $(x', y')$ reprezentują skorygowane współrzędne w układzie mapy, a $\lambda$ stanowi współczynnik skali układu współrzędnych jednorodnych.

W systemie proces ten jest aplikowany do punktu reprezentującego środek dolnej krawędzi ramki
otaczającej, stanowiącego estymowany punkt styku pieszego z podłożem. Do realizacji przekształcenia
w czasie rzeczywistym wykorzystano funkcję \texttt{cv2.perspectiveTransform} z biblioteki OpenCV.
Dzięki natywnej optymalizacji biblioteki i integracji z pakietem Apple Accelerate na architekturze
układu M5, operacje te ulegają wektoryzacji i są wykonywane w ułamkach milisekund, nie obciążając
głównego wątku analitycznego.

> **Wskazówka graficzna 1 (Rzutowanie IPM):** W tym miejscu rekomendowane jest wstawienie grafiki składającej się z dwóch obok siebie położonych obrazów. Po lewej widok z pojedynczej kamery z naniesionym wielokątem (np. obrysowanym obszarem skrzyżowania). Po prawej stronie idealnie płaska mapa z góry (BEV), na którą ten sam wielokąt jest przekształcony w prostokąt. Pomiędzy obrazami strzałka z podpisem "Homografia $\mathbf{H}$".

\section{Ograniczenia rzutowania i identyfikacja problemu okluzji}

Wykorzystanie statycznej macierzy homografii dostarcza wysoce precyzyjnych wyników pod warunkiem,
że estymowany punkt referencyjny faktycznie leży na płaszczyźnie gruntu. Zjawiska występujące w
zatłoczonym środowisku miejskim drastycznie naruszają to założenie, prowadząc do znaczących
odchyleń w wyznaczaniu globalnej pozycji.

\subsection{Zjawisko błędu paralaksy w gęstym ruchu}
Gdy przechodnie poruszają się w zwartych grupach, dochodzi do wzajemnego przesłaniania się obiektów
(okluzji). W takiej sytuacji moduł detekcji YOLOv8 wyznacza dolną krawędź ramki otaczającej dla
obiektu zasłoniętego na wysokości ciała osoby znajdującej się bliżej obiektywu, a nie na
rzeczywistym podłożu. Rzutowanie takiego punktu za pomocą macierzy IPM prowadzi do tzw. błędu
paralaksy – sztucznego odsunięcia obiektu na mapie w kierunku osi optycznej kamery. Konsekwencją
tego błędu jest drastyczne zafałszowanie odległości euklidesowych pomiędzy pieszymi w module
\texttt{GlobalTracker}, co prowadzi do błędnego działania algorytmu asocjacji i niezamierzonego
łączenia różnych tożsamości. Wymusiło to konieczność implementacji niezależnej weryfikacji
zajętości przestrzeni opartej na fizycznym ruchu.

\section{Modelowanie tła i ekstrakcja fizycznego ruchu}

Aby zweryfikować, czy obszar zadeklarowany przez moduł głębokiego uczenia jest faktycznie zajęty
przez pieszego, zaimplementowano proces wyodrębniania ruchomego pierwszego planu. Wykorzystano do
tego klasyczne techniki analizy tła (\textit{ang. Background Subtraction} -- BGS).

\subsection{Inicjalizacja algorytmu mieszanin gaussowskich (MOG2)}

Do detekcji zmian w pikselach wykorzystano algorytm \texttt{cv2.createBackgroundSubtractorMOG2}. W
odróżnieniu od podejść opartych na gęstych sieciach głębokich, MOG2 charakteryzuje się bardzo
niskim narzutem na pamięć operacyjną, modelując każdy piksel jako adaptacyjną mieszaninę rozkładów
normalnych. Mechanizm automatycznej eliminacji komponentów gaussowskich (oparty na ujemnym
współczynniku rozkładu a priori Dirichleta) skutecznie uodparnia maskę binarnego ruchu na wolno
zmieniające się warunki oświetleniowe oraz dynamiczne tło, takie jak drgająca na wietrze roślinność

\subsection{Przetwarzanie morfologiczne masek binarnego ruchu}

Surowy obraz wygenerowany przez moduł MOG2 jest podatny na występowanie szumu pomiarowego w postaci
izolowanych, pojedynczych pikseli. Z tego względu zaimplementowano krok przetwarzania
morfologicznego. Maska ruchu poddawana jest operacji otwarcia morfologicznego
(\texttt{cv2.morphologyEx} z flagą \texttt{cv2.MORPH\_OPEN}), co matematycznie odpowiada operacji
erozji, po której następuje dylatacja. Zastosowanie odpowiednio dobranego elementu strukturalnego
pozwala na usunięcie szumu oraz wygładzenie krawędzi kinetycznych, generując czysty, rzeczywisty
obraz ruchu fizycznego na scenie.

\begin{table}[htbp]
  \centering
  \caption{Parametry konfiguracyjne modułu analizy tła MOG2 i operacji morfologicznych. Źródło: opracowanie własne.}
  \label{tab:params_mog2}
  \renewcommand{\arraystretch}{1.3}
  \begin{tabular}{|l|c|p{8cm}|}
    \hline
    \textbf{Parametr}      & \textbf{Wartość} & \textbf{Opis wpływu na system}                                                    \\ \hline
    \texttt{history}       & 500              & Rozmiar okna czasowego (liczba klatek) wpływającego na modelowanie rozkładów tła. \\ \hline
    \texttt{varThreshold}  & 16.0             & Próg wariancji decydujący o zaklasyfikowaniu piksela do przestrzeni ruchu.        \\ \hline
    \texttt{detectShadows} & True             & Umożliwia identyfikację i ignorowanie cieni rzucanych przez pieszych na jezdnię.  \\ \hline
    Element strukturalny   & $3 \times 3$     & Macierz (kernel) stosowana do operacji morfologicznego otwarcia maski ruchu.      \\ \hline
  \end{tabular}
\end{table}

\section{Implementacja fuzji probabilistycznej mapy zajętości (POM)}

Wiedza pozyskana z detektora YOLO (oparta na głębokich cechach semantycznych) oraz maski ruchu z
MOG2 (oparta na zmianach pikseli) ulega fuzji w module Probabilistycznej Mapy Zajętości (POM).
Implementacja tego etapu w kodzie realizuje matematyczne założenia minimalizacji błędu pomiędzy
rzutem syntetycznym a obrazem rzeczywistym.

\subsection{Generowanie rzutów syntetycznych (Obraz A)}

Zgodnie z koncepcją POM, w pierwszej kolejności system tworzy syntetyczny obraz $A$. W pamięci
operacyjnej inicjowana jest pusta (czarna) macierz \texttt{numpy.zeros} odpowiadająca
rozdzielczości klatki kamery. Następnie skrypt iteruje po wszystkich ramkach otaczających
wygenerowanych przez model YOLO. W miejscu wystąpienia predykcji, na czarnej planszy rysowane są
białe, wypełnione prostokąty (\texttt{cv2.rectangle}). Stanowi to matematyczną hipotezę modułu AI o
geometrycznym rozkładzie zajętości w danej klatce.

\subsection{Mechanizm weryfikacji karania za niezgodność (Operacja XOR)}

Wygenerowany syntetyczny obraz rzutów ($A$) zestawiany jest ze skorygowaną maską fizycznego ruchu
pochodzącą z modułu MOG2 ($B$). Zgodnie z literaturą, miara pseudoodległości wykorzystywana do
weryfikacji zajętości komórek definiowana jest równaniem: $$ \Psi(B,A) = \frac{1}{\sigma} \frac{|B
    \otimes (1 - A) + (1 - B) \otimes A|}{|A|} $$ Licznik powyższego równania stanowi reprezentację
matematycznej różnicy symetrycznej (alternatywy wykluczającej). Aby zapewnić wydajność potoku,
operację tę zaimplementowano z wykorzystaniem funkcji \texttt{cv2.bitwise\_xor}.

Po nałożeniu obu masek piksele, które pokrywają się w obu obrazach (YOLO deklaruje człowieka i MOG2
rejestruje w tym miejscu ruch), są uznawane za zweryfikowane geometrycznie. Z kolei obszary
konfliktu (np. ramka YOLO "wisząca" w powietrzu na skutek błędu paralaksy, gdzie MOG2 zgłasza
statyczne tło) są identyfikowane i nakładana jest na nie kara geometryczna. Powoduje to odrzucenie
błędnie wyestymowanego styku obiektu z podłożem, korygując ostateczne mapowanie do przestrzeni BEV.

> **Wskazówka graficzna 2 (Ilustracja procesu POM - Różnica Symetryczna):** Jest to idealne miejsce na dodanie zestawienia 3 zminiaturyzowanych kadrów wideo.
> 1. Obraz "A": Czarne tło i na nim idealne białe kwadraty w miejscach detekcji YOLO.
> 2. Obraz "B": Surowy, zaszumiony, rzeczywisty biało-czarny obraz ruchu wygenerowany przez MOG2.
> 3. Wynik (Operacja XOR): Obraz prezentujący tylko błędy (piksele konfliktu), które pozwalają algorytmowi skorygować błąd rzutu.

\section{Fuzja w module asocjacji i globalna mapa BEV}

Informacje geometryczne, pomyślnie przetransformowane oraz skorygowane przez mechanizm POM,
przekazywane są do modułu głównego (\texttt{GlobalTracker}) jako argumenty do weryfikacji
tożsamości.

\subsection{Kalkulacja kosztu przestrzennego}

Skorygowane współrzędne z lotu ptaka służą bezpośrednio do obliczenia dystansu przestrzennego na
mapie MTMC. Zastosowano standardową metrykę euklidesową: $$ d(p_1, p_2) = \sqrt{(x_2 - x_1)^2 +
    (y_2 - y_1)^2} $$ Wartość ta stanowi komponent wagowy w ogólnej macierzy kosztów asocjacji
(omawianej w poprzednim rozdziale), będąc bezwzględnym weryfikatorem poprawności łączenia
tożsamości. Ustalenie fizycznego dystansu chroni system przed "teleportowaniem" tożsamości pieszego
do innej części skrzyżowania wyłącznie na podstawie zbieżności ubioru.

\subsection{Renderowanie ostatecznej, zunifikowanej przestrzeni MTMC}

Zwieńczeniem całego potoku przetwarzania jest wygenerowanie metrycznej, cyfrowej reprezentacji
obserwowanego obszaru w czasie rzeczywistym. System posiada dynamicznie alokowaną siatkę zajętości,
na której nanosi weryfikowane globalne tożsamości pieszych w postaci punktów. Każda naniesiona
trajektoria wzbogacona jest o metadane (Globalne ID oraz macierzysty węzeł kamery). Wynikowy rzut z
lotu ptaka nie tylko ułatwia wizualną inspekcję poprawności działania całego systemu kalibracyjnego
i reidentyfikacyjnego, ale przede wszystkim dostarcza znormalizowanych danych gotowych do analiz i
ewaluacji za pomocą dedykowanych metryk statystycznych, takich jak np. ewaluator TrackEval.

\chapter{Badanie i analiza efektywności systemu}\label{chap:badanie_efektywności}
\chapter{Podsumowanie i wnioski}\label{chap:podsumowanie}



\bibliografia

\acronymslist
\printglossaries
\listoffigures
\listoftables
\załączniki
\end{document}

